% ============================================================
%   VibraGuard — Rapport de Projet de Fin d'Études
%   BOUALI Youssef | Master IL | UIZ Agadir | 2025-2026
% ============================================================
\documentclass[12pt, a4paper, oneside]{report}

% ---- Encodage & langue ----
\usepackage[utf8]{inputenc}
\usepackage[T1]{fontenc}
\usepackage[french]{babel}
\usepackage{lmodern}
\usepackage{pdfpages}

% ---- Mise en page ----
\usepackage[
  top=2.5cm, bottom=2.5cm,
  left=3cm,  right=2.5cm,
  headheight=20pt
]{geometry}
\usepackage{setspace}
\onehalfspacing
\setlength{\parindent}{1em}
\setlength{\parskip}{0.4em}

% ---- Couleurs (noir et blanc uniquement) ----
\usepackage[table]{xcolor}
\definecolor{bluedark}{gray}{0}
\definecolor{bluemid}{gray}{0}
\definecolor{bluelight}{gray}{0.92}
\definecolor{accent}{gray}{0}
\definecolor{graybg}{gray}{0.96}
\definecolor{codered}{gray}{0}
\definecolor{codegray}{gray}{0.4}

% ---- Typographie ----
\usepackage{microtype}
\usepackage{amsmath, amssymb}
\emergencystretch=3em
\tolerance=3000
\hbadness=10000
\vbadness=10000
\hfuzz=200pt
\vfuzz=200pt
\showboxbreadth=0
\showboxdepth=0
\raggedbottom

% ---- En-têtes & pieds de page ----
\usepackage{fancyhdr}
\pagestyle{fancy}
\fancyhf{}

% Commande utilitaire : numéro de chapitre en chiffres romains majuscules
\newcommand{\chapRoman}{\Roman{chapter}}

\newcommand{\eng}[1]{\textit{#1}}

% \rightmark = section courante (ex. "I.3. MÉTHODOLOGIE DE GESTION DE PROJET")
% En-tête: section à gauche, "Chapitre III" à droite
\renewcommand{\chaptermark}[1]{\markboth{#1}{}}
\renewcommand{\sectionmark}[1]{\markright{\thesection.\ \MakeUppercase{#1}}}

\fancyhead[L]{\small\textit{\rightmark}}
\fancyhead[R]{\small\textit{%
  \ifnum\value{chapter}>0 Chapitre \chapRoman\else \leftmark\fi}}

% Pied de page: titre | numéro | auteur
\fancyfoot[L]{\small\textit{VibraGuard --- Rapport PFE}}
\fancyfoot[C]{\small\thepage}
\fancyfoot[R]{\small\textit{BOUALI Youssef | Master IL | 2025-2026}}

\renewcommand{\headrulewidth}{0.4pt}
\renewcommand{\footrulewidth}{0.4pt}

% Style pour les pages de début de chapitre (plain)
\fancypagestyle{plain}{%
  \fancyhf{}%
  \fancyfoot[L]{\small\textit{VibraGuard --- Rapport PFE}}%
  \fancyfoot[C]{\small\thepage}%
  \fancyfoot[R]{\small\textit{BOUALI Youssef | Master IL | 2025-2026}}%
  \renewcommand{\headrulewidth}{0pt}%
  \renewcommand{\footrulewidth}{0.4pt}%
}


\newif\ifinmainchapters
\inmainchapterstrue
\newif\ifinappendixpart
\fancyhead[R]{\small\textit{%
  \ifinmainchapters
    \ifnum\value{chapter}>0
      \ifinappendixpart Annexe \Alph{chapter}\else Chapitre \chapRoman\fi
    \else \leftmark\fi
  \else
    \leftmark
  \fi}}

% ---- Numérotation des chapitres en chiffres romains ----
\renewcommand{\thechapter}{\Roman{chapter}}
% Les sections gardent leur numérotation normale : I.1, I.2, ...
\renewcommand{\thesection}{\thechapter.\arabic{section}}
\renewcommand{\thesubsection}{\thesection.\arabic{subsection}}
\renewcommand{\thesubsubsection}{\thesubsection.\arabic{subsubsection}}

% ---- Titres de chapitres/sections ----
\usepackage{titlesec}
\titleformat{\chapter}[block]
  {\normalfont\Large\bfseries}
  {\ifinappendixpart Annexe \thechapter\else Chapitre \thechapter\fi}{0.8em}{}
  [\vspace{-0.5em}{\titlerule[1pt]}\vspace{0.5em}]
\titlespacing*{\chapter}{0pt}{10pt}{20pt}

\titleformat{\section}
  {\normalfont\large\bfseries}
  {\thesection}{0.8em}{}
\titlespacing*{\section}{0pt}{14pt}{6pt}

\titleformat{\subsection}
  {\normalfont\normalsize\bfseries}
  {\thesubsection}{0.7em}{}
\titlespacing*{\subsection}{0pt}{10pt}{4pt}

\titleformat{\subsubsection}
  {\normalfont\normalsize\itshape}
  {\thesubsubsection}{0.7em}{}
\titlespacing*{\subsubsection}{0pt}{8pt}{3pt}

% ---- Tables des matières ----
\usepackage[titles]{tocloft}
\setlength{\cftbeforechapskip}{6pt}
\renewcommand{\cftchapfont}{\bfseries}
\renewcommand{\cftsecfont}{}
\renewcommand{\cftchappagefont}{\bfseries}

% ---- Tableaux ----
\usepackage{booktabs}
\usepackage{tabularx}
\usepackage{longtable}
\usepackage{multirow}
\usepackage{array}
\usepackage{colortbl}
\renewcommand{\arraystretch}{1.4}
\newcolumntype{L}[1]{>{\raggedright\arraybackslash}p{#1}}
\newcolumntype{C}[1]{>{\centering\arraybackslash}p{#1}}
\newcolumntype{R}[1]{>{\raggedleft\arraybackslash}p{#1}}

% ---- Commandes tableaux (remplacées par alternance gris clair / blanc) ----
\newcommand{\theadrow}{\rowcolor{gray!30}}
\newcommand{\evenrow}{\rowcolor{gray!12}}
\newcommand{\oddrow}{\rowcolor{white}}

% ---- Listes ----
\usepackage{enumitem}
\setlist[itemize]{leftmargin=2em, itemsep=2pt, topsep=4pt}
\setlist[enumerate]{leftmargin=2em, itemsep=2pt, topsep=4pt}

% ---- Encadrés remplacés par des environnements simples sans couleur ----
\usepackage{mdframed}

\newenvironment{infobox}[1][]{%
  \begin{mdframed}[linewidth=0.8pt, linecolor=black, backgroundcolor=white,
    innerleftmargin=8pt, innerrightmargin=8pt, innertopmargin=6pt, innerbottommargin=6pt,
    nobreak=true]
  \ifx\relax#1\relax\else\textbf{#1}\par\smallskip\fi
}{%
  \end{mdframed}%
}

\newenvironment{resultbox}{%
  \begin{mdframed}[linewidth=0.8pt, linecolor=black, backgroundcolor=white,
    innerleftmargin=8pt, innerrightmargin=8pt, innertopmargin=6pt, innerbottommargin=6pt,
    nobreak=true]
}{%
  \end{mdframed}%
}

\newenvironment{chapterintro}{%
  \begin{mdframed}[linewidth=0.8pt, linecolor=black, backgroundcolor=gray!8,
    innerleftmargin=10pt, innerrightmargin=10pt, innertopmargin=8pt, innerbottommargin=8pt,
    nobreak=true]
}{%
  \end{mdframed}%
}

% ---- Code listings ----
\usepackage{listings}
\lstset{
  basicstyle=\ttfamily\small,
  keywordstyle=\bfseries,
  commentstyle=\itshape,
  stringstyle={},
  backgroundcolor=\color{gray!8},
  frame=single, framerule=0.4pt,
  numbers=left, numberstyle=\tiny,
  breaklines=true, breakatwhitespace=true,
  tabsize=2, showstringspaces=false,
  xleftmargin=1em, xrightmargin=0.5em
}

% ---- Graphiques ----
\usepackage{graphicx}
\usepackage{float}
\usepackage{caption}
\captionsetup{font=small, labelfont=bf, labelsep=period}
\frenchsetup{SmallCapsFigTabCaptions=false}

% ---- Tikz (diagrammes) ----
\usepackage{tikz}
\usetikzlibrary{shapes.geometric, arrows.meta, positioning, fit, backgrounds, calc}

% ---- Liens hypertexte ----
\usepackage[
  colorlinks=false,
  hidelinks,
  bookmarks=true,
  bookmarksopen=true,
  pdftitle={Rapport PFE -- VibraGuard},
  pdfauthor={BOUALI Youssef},
  pdfsubject={Maintenance Prédictive IoT IA Blockchain}
]{hyperref}

% ---- Acronymes ----
\usepackage{acronym}

% ---- Commandes personnalisées ----
\newcommand{\projectname}{\textbf{VibraGuard}}
\newcommand{\kw}[1]{\textbf{#1}}
\newcommand{\tech}[1]{\textit{#1}}

% ============================================================
\begin{document}
% ============================================================

% ---- Page de garde ----
\pagenumbering{gobble}
\includepdf[pages=1,pagecommand={\thispagestyle{empty}}]{Page_Garde_Rapport_PFE.pdf}


\pagenumbering{roman}
\setcounter{page}{1}

\newpage
\thispagestyle{empty}
\vspace*{5cm}
\begin{flushright}
  \textit{À mes parents, pour leur soutien indéfectible}\\[0.5cm]
  \textit{À mes encadrants, Pr. AIT LAHCEN Yassine et Mr. BOUKHIRA Adil,}\\
  \textit{pour leur confiance et leurs précieux conseils}\\[0.5cm]
  \textit{À tous ceux qui ont contribué, de près ou de loin,}\\
  \textit{à la réalisation de ce projet}\\[1cm]
  \textit{\textbf{Je dédie ce travail.}}
\end{flushright}
\newpage

% ---- Remerciements ----
\chapter*{Remerciements}
%\addcontentsline{toc}{chapter}{Remerciements}
\markboth{Remerciements}{Remerciements}

Au terme de ce Projet de Fin d'Études, il m'est particulièrement agréable d'exprimer ma
gratitude et mes remerciements à toutes les personnes qui ont contribué à la réussite de ce
travail.

Je tiens tout d'abord à exprimer ma profonde reconnaissance à mes encadrants,
\textbf{Pr.~AIT LAHCEN Yassine} et \textbf{Mr.~BOUKHIRA Adil}, pour leur disponibilité,
leur soutien constant, leurs conseils avisés et la confiance qu'ils m'ont accordée tout au
long de ce projet. Leurs orientations méthodologiques et leur expertise ont été des guides
précieux dans la conception et le développement de la plateforme \projectname{}.

Je remercie également l'ensemble du corps professoral du Centre d'Excellence IT de la
Faculté des Sciences d'Agadir --- Université Ibn Zohr --- pour la qualité de la formation
dispensée au sein du Master d'Excellence en Ingénierie Logicielle, qui m'a fourni les
fondements théoriques et pratiques nécessaires à la réalisation de ce projet.

Mes sincères remerciements vont également aux équipes de l'Office Chérifien des Phosphates
(\textbf{OCP}) de Safi pour avoir rendu possible l'ancrage industriel de ce projet et pour
l'ouverture vers des problématiques réelles de maintenance industrielle dans le contexte de
l'Industrie 4.0.

Je ne saurais oublier de remercier chaleureusement ma famille et mes proches pour leur
soutien moral indéfectible, leur patience et leurs encouragements tout au long de mon
parcours universitaire.

Enfin, je remercie tous mes collègues de promotion pour l'ambiance de travail stimulante
et les échanges fructueux qui ont contribué à enrichir ma réflexion.

\newpage

% ---- Résumé ----
\chapter*{Résumé}
\addcontentsline{toc}{chapter}{Résumé}

Dans le contexte de la transformation numérique de l'industrie (\textbf{Industrie 4.0}), la
maintenance des équipements industriels représente un enjeu stratégique majeur. Les
défaillances non anticipées des moteurs industriels génèrent des arrêts de production coûteux
et impactent significativement la compétitivité des entreprises industrielles, dont l'OCP Safi.

Le projet \projectname{} propose une solution complète et innovante de \kw{maintenance
prédictive} pour les moteurs industriels. En combinant l'Internet des Objets (\kw{IoT}) pour
la collecte de données vibratoires en temps réel, une architecture \kw{Big Data} (Apache Kafka
et Apache Spark) pour le traitement de flux massifs, et des algorithmes d'\kw{Intelligence
Artificielle} (Random Forest) pour la détection d'anomalies, \projectname{}
permet d'anticiper les pannes avant leur survenue.

La plateforme est construite sur une architecture \kw{microservices} robuste développée avec Spring Boot, exposant des APIs REST, et s'appuie sur une interface utilisateur moderne développée en React. L'intégration de la technologie \kw{Blockchain} (Ethereum/Hardhat) pour les ordres de travail et du stockage décentralisé \kw{IPFS} pour les rapports garantit l'immutabilité et la traçabilité.

\vspace{0.5em}
\noindent\textbf{Mots-clés :} Industrie 4.0, Maintenance Prédictive, Machine Learning, Random
Forest, IoT, MQTT, Apache Kafka, Apache Spark, Spring Boot, Microservices, Blockchain,
Ethereum, IPFS, React, GCP, Kubernetes.

\newpage


\chapter*{Abstract}
\addcontentsline{toc}{chapter}{Abstract}


\vspace{0.5em}
In the context of Industry 4.0's digital transformation, the maintenance of
industrial equipment represents a major strategic challenge: unplanned motor
failures generate costly production downtime and impact the competitiveness of
industrial groups such as OCP, the world's leading exporter of phosphate-based
products.

This Final-Year Project (PFE) presents \projectname{}, a comprehensive predictive
maintenance platform designed for the industrial motors of the OCP Safi site. The
solution combines IoT sensor nodes (ESP32 with DS18B20, ACS712, and
piezoelectric/ADS1115 sensors) for real-time vibration, thermal, and electrical
monitoring; a Big Data pipeline (MQTT, Apache Kafka, Apache Spark Streaming); and a Random Forest machine-learning model achieving 92.0\% accuracy in
classifying motor health states (normal operation and fault types).

The platform is built on a robust \kw{Spring Boot} \kw{microservices} architecture  exposing \kw{REST APIs}, with a modern \kw{React/TypeScript} frontend providing real-time dashboards, work-order management, and KPI visualization (MTBF, MTTR,  availability). 

\kw{Blockchain} technology (Ethereum/Hardhat smart contracts) for work orders, combined with decentralized \kw{IPFS} storage for maintenance reports, ensures  immutability and traceability — achieving  \textbf{100\%} certified events with zero failed transactions during load testing.

The entire system is containerized with \kw{Docker} and deployed on  \kw{Kubernetes} (GCP), supported by automated \kw{CI/CD} pipelines.

All measured performance criteria (accuracy, latency, throughput, availability,
test coverage, and blockchain traceability) met or exceeded their target objectives,
validating the feasibility of an integrated IoT--Big Data--AI--Blockchain approach
to predictive maintenance in a real industrial context.

\vspace{0.5em}
\noindent\textbf{Keywords:} Industry 4.0, Predictive Maintenance, Machine Learning,
Random Forest, IoT, MQTT, Apache Kafka, Apache Spark, Microservices, Blockchain,
Ethereum, IPFS, React, GCP, Kubernetes.

\newpage

% ---- Tables ----
\tableofcontents
\newpage
\listoftables
\addcontentsline{toc}{chapter}{Liste des tableaux}
\newpage
\listoffigures
\addcontentsline{toc}{chapter}{Liste des figures}
\newpage

% ---- Table des abréviations ----
\chapter*{Table des Abréviations}
\addcontentsline{toc}{chapter}{Table des Abréviations}
\markboth{Table des Abréviations}{Table des Abréviations}

\begin{center}
\scriptsize
\begin{tabularx}{\textwidth}{L{3cm} X}
  \toprule
  \rowcolor{gray!30}
  \textbf{Abréviation} & \textbf{Signification} \\
  \midrule
  \rowcolor{white}       API       & Application Programming Interface \\
  \rowcolor{gray!12}BI        & Business Intelligence \\
  \rowcolor{white}       CID       & Content Identifier (IPFS) \\
  \rowcolor{gray!12}CI/CD     & Continuous Integration / Continuous Deployment \\
  \rowcolor{white}       DAP       & Diammonium Phosphate (Phosphate Diammonique) \\
  \rowcolor{gray!12}ETL       & Extract, Transform, Load \\
  \rowcolor{white}       GCP       & Google Cloud Platform \\
  \rowcolor{gray!12}IIoT      & Industrial Internet of Things \\
  \rowcolor{white}       IPFS      & InterPlanetary File System \\
  \rowcolor{gray!12}ITIL      & Information Technology Infrastructure Library \\
  \rowcolor{white}       IoT       & Internet of Things \\
  \rowcolor{gray!12}JWT       & JSON Web Token \\
  \rowcolor{white}       KPI       & Key Performance Indicator \\
  \rowcolor{gray!12}MAP       & Monoammonium Phosphate (Phosphate Monoammonique) \\
  \rowcolor{white}       ML        & Machine Learning \\
  \rowcolor{gray!12}MTBF      & Mean Time Between Failures \\
  \rowcolor{white}       MTTR      & Mean Time To Repair \\
  \rowcolor{gray!12}MQTT      & Message Queuing Telemetry Transport \\
  \rowcolor{white}       MongoDB   & Base de données orientée documents (NoSQL) \\
  \rowcolor{gray!12}OASIS     & Organization for the Advancement of Structured Information Standards \\
  \rowcolor{white}       OCP       & Office Chérifien des Phosphates \\
  \rowcolor{gray!12}OT        & Ordre de Travail \\
  \rowcolor{white}       ONCF      & Office National des Chemins de Fer \\
  \rowcolor{gray!12}PdM       & Predictive Maintenance (Maintenance Prédictive) \\
  \rowcolor{white}       PFE       & Projet de Fin d'Études \\
  \rowcolor{gray!12}POC       & Proof of Concept \\
  \rowcolor{white}       RBAC      & Role-Based Access Control \\
  \rowcolor{gray!12}REST      & Representational State Transfer \\
  \rowcolor{white}       SLA       & Service Level Agreement (Accord de Niveau de Service) \\
  \rowcolor{gray!12}SPA       & Single Page Application \\
  \rowcolor{gray!12}TSP       & Triple Superphosphate \\
  \rowcolor{white}       UM6P      & Université Mohammed VI Polytechnique \\
\bottomrule
\end{tabularx}
\end{center}

\newpage

% ============================================================
\pagenumbering{arabic}
\setcounter{page}{1}
% ============================================================

% ============================
\chapter*{Introduction Générale}
\addcontentsline{toc}{chapter}{Introduction Générale}
\markboth{Introduction Générale}{Introduction Générale}
% ============================

\setcounter{table}{0}

Ce chapitre introductif présente le cadre général du projet \projectname{}, le contexte industriel de l'OCP Safi, la problématique identifiée, les hypothèses de travail, la méthodologie adoptée ainsi que la structure du présent rapport.

% -----------------------------------------------
\section*{1. Contexte et Hypothèse de Travail}
% -----------------------------------------------

L'ère de l'Industrie 4.0 est caractérisée par la convergence des technologies numériques avec
les systèmes de production industrielle. Dans ce nouveau paradigme, les usines intelligentes
(\eng{Smart Factories}) s'appuient sur des réseaux de capteurs connectés, des plateformes
d'analyse de données massives et des systèmes d'intelligence artificielle pour optimiser leurs
processus et réduire leurs coûts opérationnels.

\textbf{Hypothèse de travail :} il est possible de réduire significativement les arrêts de
production non planifiés de l'OCP Safi en déployant une plateforme de maintenance prédictive
basée sur l'analyse vibratoire en temps réel, les algorithmes de Machine Learning et la
certification Blockchain des données de maintenance --- à condition de disposer d'une
architecture logicielle scalable et d'un pipeline de données fiable.

% -----------------------------------------------
\section*{2. Présentation de l'OCP}
% -----------------------------------------------

\subsection*{2.1 Définition et Historique}

L'\textbf{Office Chérifien des Phosphates (OCP)} est un groupe industriel marocain créé le
\textbf{7 août 1920} par décret royal sous le Protectorat français, à l'initiative du Général
Lyautey, afin de réserver à l'État marocain l'exploitation de ses ressources phosphatières
exceptionnelles \cite{ocpgroup}. La production commerciale démarre le
1\textsuperscript{er} mars 1921 avec l'ouverture de la première mine à Khouribga, et les
premières exportations depuis le port de Casablanca ont lieu en juillet de la même année.

\subsection*{2.3 Présentation du Département IT — OCP Safi}

Le \textbf{Département IT (Direction des Systèmes d'Information)} d'OCP Safi constitue
la colonne vertébrale numérique du complexe industriel. Il est rattaché à la
\textbf{Direction Support \& DSI} du Groupe OCP et assure la conception, le déploiement
et la maintenance de l'ensemble des systèmes d'information et des infrastructures
technologiques du site de Safi.

\subsubsection*{Missions du Département IT}

\begin{enumerate}
  \item \textbf{Gestion de l'infrastructure informatique :} administration des serveurs,
  des réseaux locaux (LAN/WAN), des équipements de télécommunication et des systèmes de
  stockage de données garantissant la continuité des opérations industrielles 24h/24 et
  7j/7.

  \item \textbf{Administration des systèmes d'information :} pilotage des applications
  métier (ERP, GMAO, MES), gestion des licences logicielles et coordination des mises à
  jour en concertation avec les directions opérationnelles.

  \item \textbf{Sécurité informatique :} mise en œuvre de la politique de sécurité du
  Groupe OCP sur le site de Safi, incluant la gestion des accès (RBAC, Active Directory),
  la supervision des intrusions (SIEM), les sauvegardes et les plans de reprise après
  sinistre (PRA/PCA).

  \item \textbf{Support utilisateurs :} assistance technique de proximité pour l'ensemble
  des collaborateurs du site, gestion des incidents de niveau 1 à 3 selon le référentiel
  \textbf{ITIL} \cite{itil2019}.

  \item \textbf{Innovation et transformation digitale :} accompagnement des projets de
  digitalisation des processus industriels, intégration de nouvelles technologies (IoT,
  IA, Cloud) dans le cadre du programme \textit{OCP Digital Transformation}.
\end{enumerate}

\subsubsection*{Rôle dans le cadre du projet VibraGuard}

Dans le contexte du présent Projet de Fin d'Études, le Département IT d'OCP Safi a joué
un rôle central à plusieurs niveaux :

\begin{itemize}
  \item \textbf{Ancrage industriel :} fourniture du contexte métier, des historiques de
  maintenance et des contraintes techniques liées à l'infrastructure existante du site.
  \item \textbf{Validation des besoins :} participation à la définition des exigences
  fonctionnelles et non-fonctionnelles de la plateforme \projectname{}, notamment les
  exigences de sécurité (authentification JWT, RBAC) et de disponibilité ($> 99{,}5\%$).
  \item \textbf{Encadrement terrain :} mise à disposition de l'expertise des équipes
  techniques pour la compréhension des équipements industriels (moteurs, capteurs) et
  des processus de maintenance en vigueur à l'usine de Safi.
  \item \textbf{Perspective de déploiement :} validation de l'adéquation de l'architecture
  microservices Kubernetes de \projectname{} avec les infrastructures cloud et réseau du
  Groupe OCP, en vue d'un déploiement pilote sur les moteurs critiques du site.
\end{itemize}

\vspace{0.5em}

% -----------------------------------------------
\section*{3. Problématique}
% -----------------------------------------------

\begin{infobox}[Problématique du projet]
\textbf{Comment concevoir et déployer une plateforme de maintenance prédictive intelligente
qui permet à l'OCP Safi de surveiller en temps réel l'état de santé de ses moteurs industriels
et d'anticiper leurs défaillances, tout en garantissant la traçabilité et l'intégrité des
données de maintenance via la Blockchain ?}
\end{infobox}

\subsubsection*{Situation d'insatisfaction et éléments déclencheurs}

L'OCP Safi souffre actuellement de plusieurs dysfonctionnements dans sa gestion de la
maintenance industrielle qui constituent les éléments déclencheurs de ce projet :

\begin{itemize}
  \item L'usine de Safi \textbf{ne dispose pas d'un système de surveillance en temps réel}
  des vibrations et de la thermique de ses moteurs industriels, rendant toute maintenance
  prédictive impossible dans l'état actuel.
  \item Les arrêts non planifiés causés par les défaillances de moteurs \textbf{génèrent des
  pertes de production significatives} et des coûts de maintenance corrective élevés.
  \item Les données de maintenance sont \textbf{dispersées dans des supports hétérogènes}
  (tableurs, papier, logiciels silos) sans traçabilité centralisée ni audit trail fiable.
  \item L'\textbf{absence d'outils analytiques} ne permet pas d'exploiter les volumes
  massifs de données générées par les équipements pour en extraire de la valeur prédictive.
  \item La maintenance préventive systématique actuelle \textbf{entraîne des interventions
  prématurées coûteuses} sur des équipements encore opérationnels.
\end{itemize}

% -----------------------------------------------
\section*{4. Méthodologie --- ITIL}
% -----------------------------------------------

La réalisation du projet \projectname{} s'appuie sur la méthode \textbf{ITIL (Information
Technology Infrastructure Library)}, référentiel de bonnes pratiques pour la gestion des
services informatiques. ITIL structure notre démarche autour du cycle de vie du service :

\begin{enumerate}
  \item \textbf{Stratégie de service :} identification des besoins métier de l'OCP et
  définition du périmètre fonctionnel de la plateforme.
  \item \textbf{Conception de service :} spécification de l'architecture microservices,
  des SLA (disponibilité $>$ 99,5\%, latence $<$ 2s) et des exigences de sécurité.
  \item \textbf{Transition de service :} développement itératif, tests d'intégration (couverture
  $>$ 70\%), packaging et déploiement Kubernetes.
  \item \textbf{Exploitation de service :} monitoring Prometheus/Grafana, gestion des incidents
  via alerting automatique, support aux utilisateurs.
  \item \textbf{Amélioration continue :} réentraînement périodique des modèles ML, collecte
  des retours utilisateurs et mise à jour des dashboard BI.
\end{enumerate}

% -----------------------------------------------
\section*{5. Benchmarking}
% -----------------------------------------------

Avant de concevoir \projectname{}, une étude comparative (\textbf{benchmarking}) des
solutions de maintenance prédictive disponibles sur le marché a été conduite, selon quatre
critères : intégration IoT temps réel, capacités IA prédictive, Blockchain/traçabilité, et
coût d'implémentation.

\begin{table}[H]
  \centering
  \scriptsize
  \begin{tabularx}{\textwidth}{L{2.2cm} C{2cm} C{2cm} C{2cm} C{2.5cm} L{2cm}}
    \toprule
    \rowcolor{gray!30}
    \textbf{Solution} &
    \textbf{IoT TR} &
    \textbf{IA Préd.} &
    \textbf{Blockchain} &
    \textbf{Coût} &
    \textbf{Verdict} \\
    \midrule
    \rowcolor{white}       IBM Maximo      & Partiel & Moyen  & Non    & Très élevé & Écarté \\
    \rowcolor{gray!12}Siemens MindSphere & Oui  & Moyen  & Non    & Élevé      & Écarté \\
    \rowcolor{white}       SAP PM          & Non     & Faible & Non    & Très élevé & Écarté \\
    \rowcolor{gray!12}Open Source Mix & Partiel & Élevé  & Partiel& Faible     & Limites intégration \\
    \rowcolor{gray!18}   \textbf{\projectname{}} & \textbf{Oui} & \textbf{Élevé} & \textbf{Oui} & \textbf{Modéré} & \textbf{Retenu} \\
    \bottomrule
  \end{tabularx}
  \caption{Benchmarking des solutions de maintenance prédictive}
\end{table}

\textit{Légende : TR = Temps Réel, IA Préd. = IA Prédictive.}

Le benchmarking confirme qu'aucune solution disponible ne couvre nativement l'ensemble des
besoins de l'OCP Safi. \projectname{} est donc développé sur mesure, combinant des
technologies open source éprouvées dans une architecture intégrée et cohérente.

% -----------------------------------------------
\section*{6. Traitement des Données et Fiabilité}
% -----------------------------------------------

La qualité et la fiabilité des données constituent un enjeu fondamental pour un système de
maintenance prédictive. \projectname{} adopte une stratégie de traitement à plusieurs niveaux :

\begin{itemize}
  \item \textbf{Sources de données :} données réelles collectées par les capteurs ESP32
  (DS18B20, ACS712, Piézo+ADS1115) déployés sur les moteurs industriels de l'OCP Safi.
  \item \textbf{Pipeline de nettoyage robuste :} imputation des valeurs manquantes (médiane),
  filtrage des outliers ($\pm 3\sigma$), normalisation StandardScaler.
  \item \textbf{Split stratifié 80/20 :} pour garantir la représentativité des évaluations.
\end{itemize}

% -----------------------------------------------
\section*{7. Problèmes Rencontrés et Solutions}
% -----------------------------------------------

\begin{table}[H]
  \centering
  \begin{tabularx}{\textwidth}{L{4cm} L{4cm} L{4cm}}
    \toprule
    \rowcolor{gray!30}
    \textbf{Problème} &
    \textbf{Impact} &
    \textbf{Solution} \\
    \midrule
    \rowcolor{gray!12}
    Problèmes de liaison MQTT - Kafka & Perte de messages et instabilité du flux & Implémentation de retries avec backoff exponentiel côté ESP32, et configuration de connecteurs Kafka robustes \\
	\rowcolor{white}
    Latence Kafka-Spark $>$ 2s initialement & SLA non respecté & Optimisation micro-batch (fenêtre 2s), tuning Spark \\
    \rowcolor{gray!12}
    Coût des transactions Ethereum mainnet & Budget dépassé & Migration vers réseau Hardhat privé + testnet \\
    \rowcolor{white}
    Erreurs de déploiement et configuration Hardhat & Blocage du développement Blockchain & Mise à jour des dépendances, utilisation de Hardhat plugins, et création de scripts de déploiement automatisés \\
	\rowcolor{gray!12}
    Complexité intégration IPFS/Spring Boot & Délai de développement & Bibliothèque IPFS HTTP API + conteneurisation Kubernetes \\
    \bottomrule
  \end{tabularx}
  \caption{Principaux problèmes rencontrés et solutions apportées}
\end{table}

% -----------------------------------------------
\section*{8. Plan du Projet}
% -----------------------------------------------

\projectname{} est structuré selon une approche itérative et incrémentale en \textbf{4 mois}. Chaque mois correspond à un jalon livrable :

\begin{enumerate}
  \item \textbf{Mois 1} --- Analyse \& Conception : étude de l'existant, spécification des
  besoins, conception de l'architecture microservices, mise en place de l'environnement
  DevOps (CI/CD, Git), POC IoT+MQTT.
  \item \textbf{Mois 2} --- Développement du c\oe{}ur : services IoT et Big Data, intégration
  Kafka+Spark, développement des modèles de Machine Learning.
  \item \textbf{Mois 3} --- Blockchain \& Interfaces : développement des Smart Contracts,
  intégration IPFS, interfaces web (React) et dashboard BI.
  \item \textbf{Mois 4} --- Tests, Déploiement \& Soutenance : tests d'intégration et de
  charge, déploiement Kubernetes sur GCP, documentation technique, préparation de la
  soutenance.
\end{enumerate}

% -----------------------------------------------
\section*{8.1. Ma Contribution au Projet}
% -----------------------------------------------

Dans le cadre de ce Projet de Fin d'Études réalisé en collaboration avec l'OCP Safi,
ma contribution couvre l'intégralité du cycle de développement logiciel, de la conception
à la mise en production. Plus précisément, j'ai :

\begin{itemize}
  \item Analysé les besoins métier de l'OCP Safi et formalisé le cahier des charges technique.
  \item Conçu l'architecture microservices de \projectname{} et les diagrammes UML associés
  (cas d'utilisation, classes, séquence).
  \item Développé le pipeline de données IoT--Big Data (MQTT, Kafka, Spark Streaming).
  \item Entraîné et évalué le modèle de Machine Learning (Random Forest, 92.0\% de précision).
  \item Implémenté les microservices backend (Spring Boot) et le frontend web (React).
  \item Développé et déployé les Smart Contracts Solidity sur réseau Hardhat.
  \item Configuré l'infrastructure DevOps (Kubernetes, GCP, CI/CD GitHub Actions).
  \item Rédigé la documentation technique et le présent rapport.
\end{itemize}

% -----------------------------------------------
\section*{8.2. Méthode de Travail}
% -----------------------------------------------

Le projet a été conduit selon la méthodologie \textbf{Agile}. La communication avec les encadrants s'est organisée
comme suit :

\begin{itemize}
  \item \textbf{Réunions de suivi} avec \textbf{Pr.~AIT LAHCEN Yassine} et
  \textbf{Mr.~BOUKHIRA Adil} : présentation de l'avancement, validation des choix
  architecturaux, levée des points bloquants.
  \item \textbf{Revues régulières} : démonstration des fonctionnalités livrées et ajustement du backlog Jira.
  \item \textbf{Visites terrain} à l'OCP Safi : observations des équipements industriels,
  collecte des données de maintenance et validation des hypothèses métier.
\end{itemize}

Cette organisation a permis de maintenir un rythme de développement soutenu, de valider
en continu les orientations du projet et d'assurer une cohérence permanente entre les
besoins métier de l'OCP Safi et les choix techniques de \projectname{}.

% -----------------------------------------------
\section*{9. Présentation du Projet VibraGuard}
% -----------------------------------------------

Le projet \projectname{} répond à ces problématiques en proposant une plateforme décisionnelle
intelligente, centralisée et intégrée. Il s'articule autour de cinq piliers technologiques :

\begin{enumerate}
  \item L'\textbf{acquisition IoT} en temps réel des données vibratoires, thermiques et
  électriques des moteurs via le protocole \tech{MQTT}.
  \item Le \textbf{traitement Big Data} des flux de données grâce à Apache Kafka et Spark Streaming.
  \item L'\textbf{intelligence artificielle} pour la détection automatique d'anomalies
  (Random Forest) et la classification des défauts moteurs.
  \item La \textbf{Blockchain} pour garantir la traçabilité et la
  certification des événements de maintenance, ainsi que \textbf{l'IPFS} pour garantir l'immutabilité.
  \item Des \textbf{interfaces utilisateurs} modernes (application web React).
\end{enumerate}

% -----------------------------------------------
\section*{10. Organisation du Rapport}
% -----------------------------------------------

Le présent rapport est structuré en \textbf{cinq chapitres} :

\begin{itemize}
  \item \textbf{Chapitre I --- Contexte et État de l'Art :} présente l'évolution de la
  maintenance industrielle, les technologies de l'Industrie 4.0 et une analyse comparative
  des solutions existantes.
  \item \textbf{Chapitre II --- Analyse et Spécification des Besoins :} formalise les
  exigences fonctionnelles et non fonctionnelles, les acteurs et les cas d'utilisation.
  \item \textbf{Chapitre III --- Architecture et Choix Technologiques :} décrit
  l'architecture microservices, le flux de données et la stack technique retenue.
  \item \textbf{Chapitre IV --- Implémentation du Pipeline de Données et de l'IA :}
  présente le pipeline de prétraitement et les résultats du modèle Random Forest.
  \item \textbf{Chapitre V --- Développement du Backend, Frontend et Blockchain :}
  détaille les composants applicatifs réalisés et les résultats de déploiement.
\end{itemize}

% ============================
\chapter{Contexte et État de l'Art}
% ============================

Ce chapitre présente l'évolution de la maintenance industrielle, les technologies de l'Industrie 4.0 qui constituent le socle de \projectname{}, et une analyse comparative des solutions existantes.

\section{L'Industrie 4.0 et la Transformation Numérique Industrielle}

L'Industrie 4.0, ou quatrième révolution industrielle, désigne la transformation profonde des
processus de production par l'intégration des technologies numériques \cite{schwab2016}.
Initiée en Allemagne dès 2011 à travers le programme gouvernemental \textit{Industrie 4.0},
cette révolution repose sur plusieurs piliers fondamentaux : l'Internet des Objets Industriel
(IIoT), les systèmes cyber-physiques (CPS), l'intelligence artificielle, le Big Data, le Cloud
Computing et la robotique avancée \cite{hermann2016}.

Dans ce contexte, les usines intelligentes (\textit{Smart Factories}) se caractérisent par leur
capacité à collecter, analyser et exploiter en temps réel des volumes massifs de données provenant
de leurs équipements pour optimiser la production, anticiper les défaillances et adapter
dynamiquement les processus industriels \cite{hermann2016, schwab2016}.

\section{L'Évolution de la Maintenance Industrielle}

La maintenance industrielle a connu une évolution progressive, passant de pratiques réactives à
des approches de plus en plus prédictives \cite{mobley2002}. On distingue classiquement quatre
générations, résumées dans le tableau ci-après :

\begin{table}[H]
  \centering
  \begin{tabularx}{\textwidth}{C{1.5cm} L{3cm} L{4cm} L{4.5cm}}
    \toprule
    \rowcolor{gray!30}
    \textbf{Gén.} &
    \textbf{Stratégie} &
    \textbf{Caractéristiques} &
    \textbf{Limites} \\
    \midrule
    \rowcolor{white}
    1\textsuperscript{ère} & Corrective (Run-to-failure) & Intervention après panne & Arrêts imprévus, coûts élevés \\
    \rowcolor{gray!12}
    2\textsuperscript{ème} & Préventive systématique & Maintenance à intervalles fixes & Remplacement prématuré de pièces \\
    \rowcolor{white}
    3\textsuperscript{ème} & Préventive conditionnelle & Basée sur l'état réel mesuré & Requiert des seuils prédéfinis \\
    \rowcolor{gray!12}
    4\textsuperscript{ème} & \textbf{Prédictive (PdM)} & \textbf{Prédiction par IA sur données IoT} & \textbf{Complexité technique élevée} \\
    \bottomrule
  \end{tabularx}
  \caption{Évolution des stratégies de maintenance industrielle}
\end{table}

La maintenance prédictive (PdM) représente l'état de l'art en matière de gestion des actifs
industriels \cite{benhanifia2025}. Selon Mobley \cite{mobley2002} et le U.S. Department of Energy \cite{doe2010}, elle permet de réduire les
coûts de maintenance de 25 à 30\% et les arrêts non planifiés de 70 à 75\% par rapport à la
maintenance corrective, en déclenchant les interventions uniquement lorsque l'état réel
de l'équipement le justifie.

\section{L'Analyse Vibratoire comme Indicateur de Santé}

L'analyse vibratoire est l'une des techniques les plus efficaces pour la surveillance des
machines tournantes \cite{randall2011}. Les vibrations d'un moteur sont directement liées
à son état mécanique et peuvent révéler de nombreux types de défauts :

\begin{itemize}
  \item \textbf{Déséquilibre de masse :} augmentation de la composante vibratoire à la fréquence
  de rotation (1X) \cite{randall2011}.
  \item \textbf{Désalignement d'arbre :} présence d'harmoniques à 2X et 3X de la fréquence de rotation \cite{dewell1984}.
  \item \textbf{Défauts de roulements :} fréquences caractéristiques des bagues interne et externe
  (BPFI, BPFO), calculables à partir des paramètres géométriques du roulement \cite{randallandantoni2011}.
  \item \textbf{Surchauffe et surcharge électrique :} corrélation entre signaux de température,
  de courant et d'état mécanique, permettant une détection multi-physique \cite{lee2014}.
\end{itemize}

Randall \cite{randall2011} souligne que la combinaison de plusieurs indicateurs spectraux
(RMS, kurtosis, facteur de crête) améliore significativement la précision de détection par
rapport à l'utilisation d'un indicateur unique.

\section{L'IoT et le Big Data dans l'Industrie}

Les capteurs IoT modernes peuvent échantillonner à des fréquences élevées (plusieurs kHz),
générant des flux de données massifs difficiles à traiter avec des architectures traditionnelles.
Le protocole MQTT (\textit{Message Queuing Telemetry Transport}), standardisé par l'\textbf{OASIS}
(Organization for the Advancement of Structured Information Standards) \cite{mqtt}, s'est imposé
comme référence pour la communication légère entre capteurs embarqués et serveurs, grâce à son
architecture publish/subscribe à faible empreinte mémoire. Deux solutions s'imposent comme
standards industriels pour la gestion de ces flux :

\begin{itemize}
  \item \textbf{Apache Kafka \cite{kreps2011} :} plateforme de messagerie distribuée à haut débit,
  capable de traiter des millions de messages par seconde avec une latence inférieure à la
  milliseconde, conçue par LinkedIn et cédée à l'Apache Software Foundation en 2011.
  \item \textbf{Apache Spark Streaming \cite{spark} :} moteur de traitement distribué permettant
  l'analyse de flux de données en micro-batchs, avec intégration native des bibliothèques ML
  (MLlib) et compatibilité complète avec l'écosystème Hadoop/HDFS.
\end{itemize}

\section{L'Intelligence Artificielle pour la Détection d'Anomalies}

\subsection{Algorithmes Supervisés — Random Forest}

Introduit par Breiman \cite{breiman2001}, le Random Forest est un algorithme d'ensemble
qui construit un grand nombre d'arbres de décision indépendants et agrège leurs prédictions
par vote majoritaire. Il est particulièrement adapté à la détection de défauts industriels
car :

\begin{itemize}
    \item Il est robuste au surapprentissage grâce à l'agrégation de nombreux arbres de décision indépendants.
    \item Il gère efficacement les données multivariées complexes ainsi que les valeurs manquantes \cite{breiman2001, sklearn}.
    \item Il fournit une mesure d'importance des variables (feature importance), facilitant l'interprétabilité des décisions.
    \item Il présente de bonnes performances sur des ensembles de données de taille modérée, sans nécessiter de normalisation préalable des données.
\end{itemize}

\section{La Blockchain dans l'Industrie}

La technologie Blockchain, introduite conceptuellement par Nakamoto \cite{nakamoto2008}
et étendue aux Smart Contracts par Buterin \cite{buterin2014}, offre des propriétés uniques
dans un contexte industriel : immutabilité, transparence, traçabilité et décentralisation.
Dans \projectname{}, la Blockchain permet de :

\begin{itemize}
  \item Créer un registre immuable et infalsifiable de tous les événements de maintenance,
  conformément aux exigences de la norme ISO 55000 \cite{iso55000} et ISO 13374 \cite{iso13374}.
  \item Automatiser la validation des interventions via des Smart Contracts Solidity déployés
  sur un réseau Hardhat \cite{hardhat}.
  \item Certifier l'intégrité des rapports de maintenance via IPFS \cite{ipfs}, dont le hash
  CID est ancré on-chain de façon infalsifiable.
  \item Renforcer la confiance entre les acteurs de la maintenance (techniciens, prestataires,
  assureurs) grâce à l'auditabilité publique des enregistrements.
\end{itemize}

\subsubsection*{Pourquoi la Blockchain et non une solution classique ?}

Une base de données classique (Oracle DB, MongoDB) permettrait
certes de stocker l'historique des interventions, mais elle reste sous le contrôle
exclusif de l'administrateur système : un enregistrement peut techniquement être
modifié ou supprimé a posteriori, sans laisser de trace vérifiable de manière
indépendante. Dans un contexte industriel sensible comme celui de l'OCP, où les
rapports de maintenance peuvent être consultés par des auditeurs internes, des
assureurs ou des autorités de contrôle, cette absence de garantie d'intégrité
constitue une faiblesse.

La Blockchain apporte ici une valeur ajoutée spécifique sur quatre dimensions :

\begin{itemize}
  \item \textbf{Immutabilité :} une fois inscrit, un enregistrement (hash d'une
  intervention) ne peut plus être altéré sans que cette
  altération soit immédiatement détectable, contrairement à une simple ligne de
  base de données.
  \item \textbf{Traçabilité horodatée :} chaque transaction est associée à un
  numéro de bloc et un horodatage on-chain, constituant un journal d'audit
  infalsifiable de l'historique de maintenance d'un moteur.
  \item \textbf{Auditabilité indépendante :} la vérification d'intégrité
  (\texttt{verifyIntegrity()}) peut être effectuée par un tiers sans dépendre de
  la base de données applicative ni faire confiance à l'administrateur système.
  \item \textbf{Confiance distribuée :} dans un contexte où plusieurs parties
  prenantes (techniciens, prestataires externes, assureurs, auditeurs)
  interagissent autour des mêmes données de maintenance, la Blockchain élimine le
  besoin de faire confiance à une autorité centrale unique.
\end{itemize}

% -----------------------------------------------
\section{L'IPFS — InterPlanetary File System}
% -----------------------------------------------

\subsection{Définition et Principes Fondamentaux}

L'\textbf{InterPlanetary File System (IPFS)} est un protocole de communication et un système
de fichiers distribué pair-à-pair (\textit{peer-to-peer}), conçu par Protocol Labs et publié
en 2015 \cite{ipfs}. Contrairement au Web traditionnel fondé sur l'adressage par localisation
(URL pointant vers un serveur précis), IPFS repose sur le principe d'\textbf{adressage par
contenu} (\textit{content-addressed storage}) : chaque fichier est identifié par son empreinte
cryptographique unique, le \textbf{CID} (\textit{Content Identifier}), calculé via l'algorithme
SHA2-256.

\subsection{Fonctionnement d'IPFS}

Le fonctionnement d'IPFS s'articule autour de quatre mécanismes clés :

\begin{enumerate}
  \item \textbf{Découpage en blocs (\textit{chunking}) :} tout fichier ajouté au réseau IPFS
  est découpé en blocs de données de taille fixe. Chaque bloc reçoit un identifiant
  cryptographique unique (CID), garantissant l'intégrité du contenu.
  \item \textbf{Arbre de Merkle (\textit{Merkle DAG}) :} les blocs sont organisés en un graphe
  orienté acyclique (DAG) de type Merkle, permettant de retrouver, de vérifier et de
  déduplicer les données de façon efficace et sécurisée.
  \item \textbf{Table de hachage distribuée (\textit{DHT}) :} un annuaire distribué entre tous
  les nœuds du réseau permet de localiser quels pairs détiennent quels blocs, sans serveur
  central.
  \item \textbf{Échange de blocs (\textit{Bitswap}) :} les nœuds échangent les blocs de données
  selon un protocole d'échange pair-à-pair, assurant la disponibilité des fichiers même en
  l'absence du nœud d'origine.
\end{enumerate}

\subsection{Propriétés et Avantages d'IPFS dans VibraGuard}

L'intégration d'IPFS dans \projectname{} repose sur ses propriétés fondamentales, particulièrement
adaptées aux exigences de traçabilité et d'intégrité des données de maintenance industrielle :

\begin{table}[H]
  \centering
  \begin{tabularx}{\textwidth}{L{3.5cm} L{8cm}}
    \toprule
    \rowcolor{gray!30}
    \textbf{Propriété IPFS} & \textbf{Apport dans \projectname{}} \\
    \midrule
    \rowcolor{white}
    Adressage par contenu & Toute modification d'un rapport génère un CID différent,
    rendant toute falsification immédiatement détectable. \\
    \rowcolor{gray!12}
    Décentralisation & Les rapports d'intervention ne dépendent pas d'un serveur central
    susceptible de tomber en panne ou d'être compromis. \\
    \rowcolor{white}
    Immutabilité & Un fichier stocké sur IPFS ne peut être modifié rétroactivement :
    son CID est calculé à partir de son contenu exact. \\
    \rowcolor{gray!12}
    Déduplication & Les blocs identiques ne sont stockés qu'une seule fois sur le réseau,
    optimisant l'espace de stockage pour les rapports similaires. \\
    \rowcolor{white}
    Compatibilité Blockchain & Le CID peut être ancré directement dans un
    Smart Contract Ethereum, combinant stockage décentralisé et certification immuable. \\
    \bottomrule
  \end{tabularx}
  \caption{Propriétés d'IPFS et leur apport dans \projectname{}}
\end{table}

\subsection{Intégration IPFS dans VibraGuard}

Dans \projectname{}, IPFS fonctionne en selon le schéma suivant :

\begin{enumerate}
  \item Lorsque l'utilisateur génère un rapport PDF.
  \item Le rapport est uploadé vers le nœud IPFS via l'API HTTP \cite{ipfs}.
  \item IPFS retourne le \textbf{CID} unique du fichier (ex. : \texttt{QmXoypizjW...}).
  \item Pour vérifier l'intégrité d'un rapport, il suffit de recalculer son CID et de le
  comparer à celui enregistré : toute divergence trahit une falsification.
\end{enumerate}

Cette architecture garantit que \textbf{100\% des rapports} sont certifiés,
traçables et infalsifiables, en conformité avec les exigences des normes ISO 55000
\cite{iso55000} et ISO 13374 \cite{iso13374}.

\section{Étude de l'Existant}

Avant de concevoir \projectname{}, une analyse de l'environnement existant à l'OCP Safi a été
conduite pour identifier les pratiques actuelles de maintenance et leurs lacunes.

\subsection{État Actuel de la Maintenance à l'OCP Safi}

L'OCP Safi gère aujourd'hui la maintenance de ses moteurs industriels selon une approche
hybride :

\begin{itemize}
  \item \textbf{Fiches papier et tableurs Excel :} les données de maintenance (dates d'intervention,
  pièces remplacées, durées d'arrêt) sont saisies manuellement, engendrant des erreurs et une
  traçabilité insuffisante.
  \item \textbf{GMAO partielle :} un outil de gestion de maintenance assistée par ordinateur
  existe mais ne couvre pas la supervision temps réel ni l'analyse prédictive.
  \item \textbf{Absence de capteurs vibratoires permanents :} les mesures vibratoires sont
  effectuées ponctuellement par des techniciens avec des analyseurs portatifs, sans historisation
  automatique.
  \item \textbf{Maintenance préventive calendaire :} les interventions sont programmées à des
  intervalles fixes (ex. toutes les 3 000 heures de fonctionnement), indépendamment de l'état
  réel des équipements.
\end{itemize}

\subsection{Étude de Cas — Défaillances Typiques des Moteurs à OCP Safi}

L'analyse des historiques de maintenance disponibles révèle que les défaillances les plus
fréquentes sur le parc de moteurs industriels de l'OCP Safi sont :

\begin{table}[H]
  \centering
  \begin{tabularx}{\textwidth}{L{4.5cm} L{4cm} L{5.5cm}}
    \toprule
    \rowcolor{gray!30}
    \textbf{Type de Défaillance} & 
    \textbf{Cause Principale} & 
    \textbf{Impact} \\
    \midrule
    \rowcolor{white} 
    Défaut de roulement & Lubrification insuffisante & Arrêt imprévu 4--12h \\
    \rowcolor{gray!12}
    Déséquilibre de masse & Usure / dépôt de matière & Vibrations excessives \\
    \rowcolor{white} 
    Surchauffe moteur & Surcharge / ventilation & Grillage bobinage \\
    \rowcolor{gray!12}
    Désalignement d'arbre & Mauvais alignement install. & Usure prématurée \\
    \rowcolor{white} 
    Surcharge électrique & Variations tension secteur & Court-circuit \\
    \bottomrule
  \end{tabularx}
  \caption{Typologies de défaillances sur le parc moteurs OCP Safi}
\end{table}

Ces données justifient la pertinence des cinq classes de défauts modélisées dans \projectname{},
en cohérence avec la taxonomie des défaillances de machines tournantes proposée par Jardine
et al. \cite{jardine2006}.

\section{Étude Comparative des Solutions Existantes}

\subsection{Critères de Comparaison}

Le benchmarking des solutions a été conduit selon cinq critères pondérés reflétant les besoins
spécifiques de l'OCP Safi. La méthodologie de sélection s'appuie sur les recommandations de
gestion des actifs industriels de la norme ISO 55000 \cite{iso55000} :

\begin{table}[H]
  \centering
  \scriptsize
  \begin{tabularx}{\textwidth}{L{2.5cm} C{1.5cm} C{1.8cm} C{1.5cm} C{1.5cm} C{1.5cm} C{1.8cm}}
    \toprule
    \rowcolor{gray!30}
    \textbf{Critère} &
    \textbf{Poids} &
    \textbf{IBM Maximo} &
    \textbf{Siemens} &
    \textbf{SAP PM} &
    \textbf{Open Source} &
    \textbf{\projectname{}} \\
    \midrule
    \rowcolor{white}       IoT Temps Réel   & 25\% & Partiel & Oui  & Non    & Partiel & \textbf{Oui} \\
    \rowcolor{gray!12}IA Prédictive    & 25\% & Moyen   & Moyen& Faible & Élevé   & \textbf{Élevé} \\
    \rowcolor{white}       Blockchain       & 20\% & Non     & Non  & Non    & Partiel & \textbf{Oui} \\
    \rowcolor{gray!12}Coût             & 15\% & Très élevé & Élevé & Très élevé & Faible & \textbf{Modéré} \\
    \rowcolor{white}       Scalabilité      & 15\% & Moyen   & Élevé& Moyen  & Faible  & \textbf{Élevé} \\
    \midrule
    \rowcolor{gray!18}   \textbf{Score global} & \textbf{100\%} & 52/100 & 61/100 & 38/100 & 55/100 & \textbf{89/100} \\
    \bottomrule
  \end{tabularx}
  \caption{Comparaison détaillée des solutions de maintenance prédictive}
\end{table}

\subsection{Étude de Choix Technologique}

Sur la base de ce benchmarking et des contraintes de l'OCP Safi, le choix de développer
\projectname{} en solution sur mesure s'impose. Les justifications clés sont :

\begin{itemize}
  \item \textbf{IBM Maximo et SAP PM :} solutions propriétaires très coûteuses en licences et
  maintenance, sans intégration native Blockchain ni capacités IA prédictive avancées.
  Le coût d'implémentation dépasse les contraintes budgétaires du projet.
  \item \textbf{Siemens MindSphere :} excellente couverture IoT mais fort \textit{vendor lock-in}
  sur l'écosystème Siemens. La couche IA prédictive reste limitée aux outils Siemens et n'intègre
  pas de Blockchain.
  \item \textbf{Solutions Open Source agrégées :} flexibles mais nécessitent une intégration
  custom complexe sans garantie de cohérence architecturale ni de support enterprise.
  \item \textbf{\projectname{} sur mesure :} permet d'intégrer nativement IoT (MQTT \cite{mqtt}/Kafka
  \cite{kafka}), IA (Random Forest/Scikit-learn \cite{breiman2001, sklearn}), Blockchain
  (Ethereum \cite{ethereum}/Hardhat \cite{hardhat}) et IPFS \cite{ipfs} dans une architecture
  microservices \cite{springboot} cohérente, à un coût de développement maîtrisé dans le cadre du PFE.
\end{itemize}

L'analyse confirme que \projectname{} se positionne précisément dans le créneau non couvert
par les solutions existantes.

\section{Conclusion}

Ce chapitre a présenté le cadre théorique et technologique du projet \projectname{}. Nous
avons retracé l'évolution de la maintenance industrielle vers les approches prédictives, analysé
les technologies clés de l'Industrie 4.0 (IoT, Big Data, IA, Blockchain, IPFS) et étudié les
solutions existantes sur le marché. Le benchmarking réalisé confirme que \projectname{} répond
à un besoin non couvert par les solutions commerciales, en combinant de façon native les
capacités IoT temps réel, l'intelligence artificielle et la traçabilité Blockchain dans une
architecture microservices adaptée au contexte de l'OCP Safi. Le chapitre suivant s'attache
à la spécification formelle des besoins et à la conception architecturale de la plateforme.

% ============================
\chapter{Analyse et Spécification des Besoins}
% ============================

Ce chapitre présente les objectifs du projet, la méthodologie Agile adoptée, l'identification des acteurs, les besoins fonctionnels et non-fonctionnels, les diagrammes UML ainsi que les technologies retenues pour la réalisation de \projectname{}.

\section{Objectifs du Projet}

Le projet \projectname{} poursuit les objectifs suivants :

\begin{enumerate}
  \item \textbf{Surveillance temps réel :} développer une plateforme de monitoring continu
  de l'état de santé des moteurs industriels de l'OCP Safi, basée sur l'analyse vibratoire, thermique et électrique.
  \item \textbf{Détection automatique d'anomalies :} mettre en place un système de
  classification automatique des défauts mécaniques avec une précision supérieure à 85\%
  (dépassé : 92.0\% atteint).
  \item \textbf{Gestion des interventions :} automatiser la création, l'assignation et le
  suivi des Ordres de Travail (OT) avec workflow intégré et notifications push.
  \item \textbf{Traçabilité immuable :} garantir via Blockchain la certification et
  l'immutabilité de 100\% des événements critiques de maintenance.
  \item \textbf{Tableau de bord décisionnel :} fournir aux responsables maintenance des KPI
  en temps réel (MTBF, MTTR, taux de disponibilité) via une interface web dédiée.
  \item \textbf{Architecture scalable :} concevoir une plateforme microservices déployée sur
  Kubernetes capable d'intégrer de nouveaux capteurs sans modification du code.
\end{enumerate}

\section{Méthodologie Adoptée — Agile avec Jira}

\projectname{} est développé selon la \textbf{méthodologie Agile}, gérée
via \textbf{Jira Software}. Cette approche itérative et incrémentale permet de livrer de la
valeur et d'adapter le périmètre aux retours des parties prenantes.

\subsection{Tableau de Planification — Diagramme de Gantt}

Le diagramme de Gantt suivant illustre la répartition des tâches sur les 4 mois du projet,
structurées en 4 phases itératives avec leurs dépendances :

\begin{figure}[H]
  \centering
  \includegraphics[width=\textwidth]{Online_Gantt_20260208.png}
  \caption{Diagramme de Gantt du projet VibraGuard (Fév.--Juin 2026)}
\end{figure}

\begin{table}[H]
  \centering
  \begin{tabularx}{\textwidth}{L{1.5cm} L{2.2cm} L{5cm} L{3.5cm}}
    \toprule
    \rowcolor{gray!30}
    \textbf{Phase} &
    \textbf{Période} &
    \textbf{Activités Principales} &
    \textbf{Livrables} \\
    \midrule
    \rowcolor{white}
    Mois 1 & Fév. 2026 & Étude existant, conception archi, setup CI/CD, POC IoT+MQTT & Cahier des charges, schéma archi \\
    \rowcolor{gray!12}
    Mois 2 & Mar. 2026 & Services IoT+Big Data, Kafka+Spark, modèles ML & Pipeline fonctionnel, modèle IA \\
    \rowcolor{white}
    Mois 3 & Avr. 2026 & Smart contracts, IPFS, interfaces React, dashboard BI & Application web, blockchain \\
    \rowcolor{gray!12}
    Mois 4 & Mai-Juin 2026 & Tests, déploiement Kubernetes, documentation, soutenance & Rapport PFE, vidéo démo \\
    \bottomrule
  \end{tabularx}
  \caption{Synthèse du planning par phase}
\end{table}

\section{Identification et Définition des Acteurs}

\subsection{Définition des Acteurs Principaux}

Un \textbf{acteur} est toute entité (humaine ou système) qui interagit directement avec la
plateforme \projectname{}. On distingue trois acteurs humains et deux acteurs systèmes :

\begin{itemize}
  \item \textbf{Technicien de Maintenance :} opérateur terrain chargé d'exécuter les ordres de
  travail générés par le système. Il consulte les alertes via l'interface web,
  intervient sur les équipements défaillants et renseigne les rapports d'intervention.

  \item \textbf{Responsable Maintenance :} superviseur qui pilote les KPI globaux du parc de
  machines via le dashboard web. Il valide les rapports, planifie les interventions en fonction
  des prédictions du système et analyse les tendances de dégradation.

  \item \textbf{Administrateur Système :} gestionnaire de la plateforme informatique. Il gère
  les comptes utilisateurs (RBAC), configure les capteurs IoT, surveille l'infrastructure et
  administre les Smart Contracts.

  \item \textbf{Capteur IoT (acteur système) :} entité autonome qui publie en continu les
  données vibratoires, thermiques et électriques vers le broker MQTT.

  \item \textbf{Blockchain / Smart Contract (acteur système) :} entité autonome qui certifie
  automatiquement les événements critiques et répond aux requêtes de vérification d'intégrité.
\end{itemize}

\begin{table}[H]
  \centering
  \begin{tabularx}{\textwidth}{L{3.5cm} C{2.5cm} C{2.5cm} C{2.5cm}}
    \toprule
    \rowcolor{gray!30}
    \textbf{Acteur} &
    \textbf{IoT} &
    \textbf{OT} &
    \textbf{Blockchain} \\
    \midrule
    \rowcolor{white}       Technicien        & Lecture  & Écriture & Lecture \\
    \rowcolor{gray!12}Resp. Maintenance & Lecture      & Écriture  & Lecture \\
    \rowcolor{white}       Administrateur    & Lecture  & Écriture & Lecture \\
    \rowcolor{gray!12}Capteur IoT       & Écriture & ---      & ---     \\
    \rowcolor{white}       Smart Contract    & ---      & Lecture  & Écriture\\
    \bottomrule
  \end{tabularx}
  \caption{Matrice acteurs / modules de VibraGuard}
\end{table}

\section{Diagrammes UML}

\subsection{Diagramme de Cas d'Utilisation}

Le diagramme de cas d'utilisation ci-dessous présente l'ensemble des fonctionnalités de
\projectname{} regroupées en six modules, avec les relations entre acteurs et cas d'usage
(\textit{include} et \textit{extend}) :

\begin{figure}[H]
  \centering
  \includegraphics[width=0.55\textwidth]{usecase2.png}
  \caption{Diagramme de Cas d'Utilisation --- Plateforme VibraGuard}
\end{figure}

Le diagramme couvre des cas d'utilisation répartis dans six modules fonctionnels : Module IoT
(UC01--UC02), Module Analyse (UC03--UC04), Module Alertes (UC05--UC06), Module Maintenance
(UC07--UC09), Module Traçabilité (UC10--UC12) et Module Administration (UC15--UC16).

\subsection{Diagramme de Classes}

Le diagramme de classes modélise les entités métier de la plateforme, leurs attributs, méthodes
et relations :

\begin{figure}[H]
  \centering
  \includegraphics[width=\textwidth]{classes2.png}
  \caption{Diagramme de Classes --- Plateforme VibraGuard}
\end{figure}

Les entités principales sont : \textbf{Moteur} (cœur du domaine), \textbf{DonneeVibration}
(données capteurs temps réel), \textbf{Anomalie} (résultat de la détection IA),
\textbf{OrdreTravail} (gestion des interventions), \textbf{BlockchainRecord} (certification
immuable), \textbf{DocumentIPFS} (stockage décentralisé), \textbf{Utilisateur} (gestion des
accès RBAC) et \textbf{Alerte} (notification temps réel).

\section{Besoins Fonctionnels}

\subsection{Module IoT et Acquisition de Données}
\begin{itemize}
  \item \textbf{BF-01 :} Le système doit recevoir en temps réel les données des capteurs IoT
  (vibration, température, courant, vitesse) via le protocole MQTT.
  \item \textbf{BF-02 :} Le système doit valider et normaliser les données entrantes avant traitement.
  \item \textbf{BF-03 :} Le système doit gérer les défaillances de connexion avec des mécanismes
  de \textit{retry} automatique.
  \item \textbf{BF-04 :} Le système doit supporter la collecte simultanée de données de 10+ capteurs.
\end{itemize}

\subsection{Module Traitement et Détection d'Anomalies}
\begin{itemize}
  \item \textbf{BF-05 :} Le système doit analyser les flux de données en temps réel via Spark Streaming.
  \item \textbf{BF-06 :} Le système doit classifier l'état du moteur (Normal, Déséquilibre,
  Désalignement, Défaut de roulement, Surcharge, Surchauffe) avec une précision $> 85\%$.
  \item \textbf{BF-07 :} Le système doit générer des alertes automatiques dès détection d'anomalie.
\end{itemize}

\subsection{Module Tableau de Bord et Visualisation}
\begin{itemize}
  \item \textbf{BF-09 :} Le système doit afficher en temps réel les KPI globaux (taux de
  disponibilité, MTBF, MTTR, nombre d'alertes actives).
  \item \textbf{BF-10 :} Le système doit présenter l'état de santé de chaque moteur via des
  indicateurs visuels intuitifs (jauges circulaires, codes couleur).
  \item \textbf{BF-11 :} Le système doit intégrer des tableaux de bord analytiques.
\end{itemize}

\subsection{Module Gestion des Ordres de Travail}
\begin{itemize}
  \item \textbf{BF-13 :} Le système doit permettre la création, l'assignation et le suivi des Ordres de Travail (OT).
  \item \textbf{BF-14 :} Le système doit notifier automatiquement les techniciens assignés.
  \item \textbf{BF-15 :} Le système doit générer des rapports d'intervention à la clôture de chaque OT.
\end{itemize}

\subsection{Module Blockchain et Traçabilité}
\begin{itemize}
  \item \textbf{BF-16 :} Le système doit enregistrer de façon immuable sur la blockchain le
  hash cryptographique de chaque intervention.
  \item \textbf{BF-17 :} Le système doit automatiser la validation via des Smart Contracts.
  \item \textbf{BF-19 :} Le système doit stocker les rapports volumineux sur IPFS et enregistrer
  leur CID sur la blockchain.
\end{itemize}

\section{Besoins Non-Fonctionnels}

\begin{table}[H]
  \centering
  \begin{tabularx}{\textwidth}{L{2.5cm} L{4.5cm} C{2.5cm}}
    \toprule
    \rowcolor{gray!30}
    \textbf{Catégorie} &
    \textbf{Exigence} &
    \textbf{Valeur Cible} \\
    \midrule
    \rowcolor{white}
    Performance & Latence traitement IoT & $< 2$\,s \\
    \rowcolor{gray!12}
    Performance & Débit de traitement & $> 10\,000$\,msg/s \\
    \rowcolor{white}
    Performance & Temps de réponse API & $< 200$\,ms (P95) \\
    \rowcolor{gray!12}
    Disponibilité & Taux de disponibilité & $> 99{,}5\%$ \\
    \rowcolor{white}
    Sécurité & Authentification & JWT \\
    \rowcolor{gray!12}
    Conformité & Standard sécurité & ISO 27001 \\
    \rowcolor{gray!12}
    Précision IA & Détection d'anomalies & $> 85\%$ (atteint : $92{,}0\%$) \\
    \bottomrule
  \end{tabularx}
  \caption{Exigences non-fonctionnelles du système}
\end{table}

\section{Technologies Adoptées}

Le choix des technologies a été guidé par quatre critères : maturité de l'écosystème,
performance, coût (open source prioritaire) et adéquation aux besoins spécifiques de
\projectname{}.

\begin{table}[H]
  \centering
  \begin{tabularx}{\textwidth}{L{2.2cm} L{3.8cm} L{5.5cm}}
    \toprule
    \rowcolor{gray!30}
    \textbf{Couche} &
    \textbf{Technologie} &
    \textbf{Justification du Choix} \\
    \midrule
    \rowcolor{white}
    Backend & Java 17 + Spring Boot 3 & Standard entreprise, écosystème Spring Cloud mature,
    microservices natifs \\
    \rowcolor{gray!12}
    Frontend & React 18 + TypeScript + Tailwind & Réactivité composants, typage statique,
    design system cohérent \\
    \rowcolor{white}
    IoT & MQTT + Eclipse Mosquitto & Protocole léger optimisé IoT, QoS configurable \\
    \rowcolor{gray!12}
    Messagerie & Apache Kafka & Haut débit ($>10M$ msg/s), persistance, rejeu possible \\
    \rowcolor{white}
    Big Data & Apache Spark & Moteur distribué, intégration MLlib native \\
    \rowcolor{gray!12}
    ML & Scikit-learn + Python 3.10 & Bibliothèques matures, Random Forest performant \\
    \rowcolor{white}
    Blockchain & Ethereum + Solidity + Hardhat & Smart contracts Turing-complets, réseau privé \\
    \rowcolor{gray!12}
    Stockage & IPFS + Oracle DB + MongoDB & Rapports décentralisés + BD relationnelle + BD vibrations \\
    \rowcolor{white}
    DevOps & Kubernetes + GitHub Actions & Orchestration conteneurs, auto-scaling, CI/CD \\
    \rowcolor{gray!12}
    Monitoring & Prometheus & Observabilité complète, alerting proactif \\
    \rowcolor{white}
    Gestion projet & Jira Software & Agile, burndown, Kanban intégrés \\
    \bottomrule
  \end{tabularx}
  \caption{Justification des choix technologiques}
\end{table}

\section{Conclusion}

Ce chapitre a défini de manière rigoureuse les besoins fonctionnels et non fonctionnels de
\projectname{}, en s'appuyant sur une identification précise des acteurs et une modélisation
UML complète. Les exigences sont formalisées en termes mesurables (précision IA $>$ 85\%,
latence $<$ 2 s, disponibilité $>$ 99,5\%) permettant une validation objective en fin de
projet. Les technologies retenues ont été justifiées sur des critères objectifs de maturité,
performance et adéquation au contexte industriel. Le chapitre suivant détaille l'architecture
technique mise en œuvre pour répondre à ces exigences.

% ============================
\chapter{Architecture et Choix Technologiques}
% ============================

Ce chapitre décrit l'architecture microservices de \projectname{}, le flux de données \textit{end-to-end}, l'infrastructure cloud GCP et la stack technique retenue.

\section{Architecture Microservices}

L'architecture de \projectname{} se décompose en \textbf{microservices indépendants},
chacun responsable d'un domaine fonctionnel précis, selon les principes de découplage
et de responsabilité unique préconisés dans le développement d'applications cloud-natives
\cite{springboot, kubernetes} :

\begin{table}[H]
  \centering
  \begin{tabularx}{\textwidth}{L{3.5cm} L{4.5cm} L{4.5cm}}
    \toprule
    \rowcolor{gray!30}
    \textbf{Composant / Service} &
    \textbf{Responsabilités} &
    \textbf{Technologies Clés} \\
    \midrule
    \rowcolor{white}
    Passerelle IoT (Bridge) & Réception MQTT et publication vers Kafka & Python, Mosquitto, Kafka \\
    \rowcolor{gray!12}
    Pipeline Streaming \& IA & Consommation Kafka, Spark Streaming, détection d'anomalies (Random Forest) & Python, PySpark, Scikit-learn \\
    \rowcolor{white}
    Backend API (Core) & API REST, gestion des moteurs, persistance des données, Gateway & Java, Spring Boot, Oracle DB, MongoDB \\
    \rowcolor{gray!12}
    Blockchain & Smart contracts, certification d'interventions, audit trail & Solidity, Hardhat, Ethers.js \\
    \rowcolor{white}
    IPFS & Stockage décentralisé des rapports d'intervention & IPFS Node (Kubo) \\
    \rowcolor{gray!12}
    Frontend (Dashboards) & Interface utilisateur, visualisation temps réel, KPI & React, TypeScript, Vite, Tailwind CSS \\
    \bottomrule
  \end{tabularx}
  \caption{Microservices et Composants de la plateforme VibraGuard}
\end{table}

\subsection{Vue d'Ensemble de l'Architecture}

La figure suivante présente le schéma global de l'architecture pipeline de \projectname{},
illustrant les cinq couches technologiques et leurs interactions :

\begin{figure}[H]
  \centering
  \includegraphics[width=\textwidth]{schema_de_projet.png}
  \caption{Architecture pipeline de \projectname{} — Vue d'ensemble des cinq couches
  (I/O Edge, Big Data, IA, Backend Microservices, Interfaces utilisateur)}
\end{figure}

Ce schéma synthétise le flux de données de bout en bout : les moteurs industriels
(Moteur 1 à N) publient leurs mesures (vibration, température, courant) via le
protocole \textbf{MQTT} vers le broker Eclipse Mosquitto, puis vers \textbf{Apache Kafka}
(10k msg/s). La couche Big Data (\textbf{Apache Spark Streaming}, micro-batchs de 2 s)
consomme ces flux et les transmet au modèle \textbf{Random Forest} (Python/Scikit-learn,
92,0\% de précision). Les résultats alimentent le backend \textbf{Spring Boot / Java 17}
via API Gateway, qui orchestre les services Motor, Oracle DB, MongoDB. Enfin, la
couche présentation expose les données en temps réel via \textbf{WebSocket}, les tableaux
de bord \textbf{BI}, et assure la traçabilité via \textbf{Blockchain/IPFS} et
l'interface \textbf{React/TypeScript}.

\section{Flux de Données End-to-End}

Le parcours d'une donnée capteur au sein de \projectname{} suit le chemin suivant,
en s'appuyant sur le protocole MQTT \cite{mqtt} pour l'acquisition et Apache Kafka
\cite{kafka} pour le transport à haut débit.

% ---- Schéma 1 : Pipeline MQTT / Kafka / Spark ----
\subsection{Pipeline IoT — MQTT, Kafka et Spark}

\begin{figure}[H]
  \centering
  \includegraphics[width=\textwidth]{MQTT_Kafka_Spark_schema.png}
  \caption{Architecture du pipeline temps réel IoT --- Intégration MQTT, Apache Kafka
  et Apache Spark Streaming dans \projectname{} \cite{mqtt, kafka, spark}}
  \label{fig:mqtt_kafka_spark}
\end{figure}

La figure~\ref{fig:mqtt_kafka_spark} illustre le pipeline de données temps réel en quatre
couches. \textbf{(1) Edge / Ingestion MQTT} : les capteurs IoT publient leurs mesures
(vibration, température, courant) via le protocole MQTT \cite{mqtt} vers le broker
Eclipse Mosquitto, qui les transmet au connecteur MQTT-Kafka. \textbf{(2) Backbone
Kafka} : les messages sont ingérés dans des topics Apache Kafka \cite{kafka,kreps2011}
dédiés par type de mesure, avec réplication et partitionnement garantissant un débit
supérieur à $10\,000$ messages/seconde. \textbf{(3) Traitement Spark} : Apache Spark
Streaming \cite{spark} consomme les partitions Kafka en micro-batchs de 2 secondes,
applique les transformations (fenêtrage, filtrage, inférence ML) et produit un DataFrame
de résultats. \textbf{(4) Données de sortie} : les résultats sont acheminés vers le
stockage froid (MongoDB), les bases temps-réel et les tableaux de bord.

% ---- Schéma 2 : Blockchain Hardhat ----
\subsection{Écosystème Blockchain — Hardhat}

\begin{figure}[H]
  \centering
  \includegraphics[width=\textwidth]{Hardhat_blockchain_schema.png}
  \caption{Écosystème de développement Blockchain Hardhat --- Configuration,
  compilation, tests et déploiement des Smart Contracts de \projectname{} \cite{hardhat}}
  \label{fig:hardhat}
\end{figure}

La figure~\ref{fig:hardhat} décrit le workflow complet de développement Blockchain
adopté dans \projectname{}, structuré en trois volets. \textbf{(1) Configuration et
schéma} (\texttt{hardhat.config.ts}) : définition du compilateur Solidity (v0.8.20),
des réseaux cibles (réseau local Hardhat en mémoire, \textit{localhost} JSON-RPC,
testnet) et des chemins sources (\texttt{contracts/}, \texttt{test/},
\texttt{artifacts/}). \textbf{(2) Système de tâches Hardhat Runner} : la CLI Hardhat
orchestre les commandes \texttt{init}, \texttt{compile} (génération ABI/bytecode via
\texttt{solc}), \texttt{test} (TypeScript/JS et tests Solidity Foundry-compatibles),
\texttt{run/deploy} et \texttt{verify}. Un réseau
Hardhat local EVM intégré (modes \textit{in-process} et \textit{standalone}, minage
automatique, traces de débogage) permet des tests sans coût de gaz. \textbf{(3)
Workflow de développement} : de l'initialisation (\texttt{npx hardhat -{}-init}) à la
vérification du contrat déployé, en passant par l'écriture des Smart Contracts Solidity
\cite{ethereum}, la génération d'artefacts, les tests, le déploiement local, puis le
déploiement sur réseau live et la vérification.

Dans \projectname{}, ce workflow a permis de déployer le contrat
\texttt{WorkOrderRegistry.sol} sur un réseau Hardhat privé, garantissant des
transactions sans frais de gaz et une latence de confirmation de 0 seconde en
environnement de développement, avant une éventuelle migration vers un testnet public.

% ---- Schéma 3 : IPFS ----
\subsection{Architecture IPFS — Stockage Décentralisé}

\begin{figure}[H]
  \centering
  \includegraphics[width=\textwidth]{IPFS_schema.png}
  \caption{Architecture du système IPFS --- Ajout de fichier côté client (Merkle DAG),
  fonctionnement réseau (DHT + Bitswap) et récupération par adressage par contenu (CID) \cite{ipfs}}
  \label{fig:ipfs_schema}
\end{figure}

La figure~\ref{fig:ipfs_schema} illustre les trois mécanismes fondamentaux d'IPFS
\cite{ipfs} exploités dans \projectname{}. \textbf{(1) Ajout d'un fichier côté client} :
lorsqu'un rapport PDF est soumis à IPFS, il est découpé en blocs (\textit{chunking}),
chaque bloc étant haché individuellement. Les blocs sont organisés en un arbre de Merkle
(\textit{Merkle DAG}), dont la racine constitue le \textbf{CID} (\textit{Content
Identifier}) --- empreinte cryptographique unique du fichier calculée via SHA2-256.
\textbf{(2) Fonctionnement réseau} : une table de hachage distribuée (\textit{DHT})
permet à tout nœud du réseau de localiser les pairs détenant les blocs recherchés, sans
serveur central. Le protocole \textit{Bitswap} assure ensuite l'échange pair-à-pair des
blocs. \textbf{(3) Adressage par contenu et récupération} : l'utilisateur soumet une
requête CID ; IPFS localise les pairs détenteurs des blocs et les télécharge en
parallèle, puis reconstruit le fichier original.

Dans \projectname{}, cette architecture garantit que tout rapport stocké
sur IPFS est \textbf{infalsifiable} : la moindre modification du contenu produirait un
CID différent, immédiatement détectable par comparaison avec le CID ancré \textit{on-chain}
dans le Smart Contract Ethereum \cite{buterin2014}.

% ---- Liste numérotée du flux end-to-end ----
\vspace{0.5em}
\noindent Le flux de données complet de \projectname{} s'articule donc selon les étapes suivantes :

\begin{enumerate}
  \item \textbf{Acquisition :} Le capteur IoT envoie ses mesures via MQTT \cite{mqtt}
  au broker Eclipse Mosquitto (figure~\ref{fig:mqtt_kafka_spark}, couche 1).
  \item \textbf{Ingestion :} Le Service IoT Bridge consomme les messages MQTT et les
  publie dans des topics Apache Kafka \cite{kafka} dédiés (figure~\ref{fig:mqtt_kafka_spark},
  couche 2).
  \item \textbf{Traitement :} Apache Spark Streaming \cite{spark} consomme les messages
  Kafka en micro-batchs de 2 secondes, applique le prétraitement et exécute le modèle
  Random Forest \cite{breiman2001,sklearn} (figure~\ref{fig:mqtt_kafka_spark}, couche 3).
  \item \textbf{Stockage :} Les données de vibration sont archivées dans MongoDB ; les
  données métier dans Oracle DB, conformément aux exigences de traçabilité ISO 13374
  \cite{iso13374}.
  \item \textbf{Certification :} Les interventions sont hachées et inscrites sur
  la Blockchain Ethereum \cite{ethereum,buterin2014} via Hardhat
  (figure~\ref{fig:hardhat}) ; les rapports PDF sont stockés sur IPFS \cite{ipfs}
  (figure~\ref{fig:ipfs_schema}).
  \item \textbf{Visualisation :} Le Frontend React \cite{react} consomme les APIs REST
  du backend Spring Boot \cite{springboot}.
\end{enumerate}


\section{Stack Technique Complète}

\begin{table}[H]
  \centering
  \begin{tabularx}{\textwidth}{L{2.5cm} L{4.5cm} L{4cm}}
    \toprule
    \rowcolor{gray!30}
    \textbf{Couche} &
    \textbf{Technologies} &
    \textbf{Justification} \\
    \midrule
    \rowcolor{white}
    Backend & Java 17, Spring Boot 3, Spring Cloud & Maturité, performance \\
    \rowcolor{gray!12}
    Frontend Web & React 18, TypeScript, Tailwind CSS & Réactivité, DX optimale \\
    \rowcolor{white}
    Blockchain & Ethereum, Solidity, Hardhat & Smart contracts robustes \\
    \rowcolor{gray!12}
    Messagerie & Apache Kafka, MQTT & Haut débit, découplage \\
    \rowcolor{white}
    Big Data & Apache Spark & Traitement distribué temps réel \\
    \rowcolor{gray!12}
    ML & Python, Scikit-learn, Pandas, NumPy & Bibliothèques ML matures \\
    \rowcolor{white}
    Base de données & Oracle DB + MongoDB & BD relationnelle + BD vibrations NoSQL \\
    \rowcolor{white}
    Orchestration & Kubernetes, Helm & Scalabilité, reproductibilité \\
    \rowcolor{gray!12}
    Monitoring & Prometheus & Observabilité complète \\
    \bottomrule
  \end{tabularx}
  \caption{Stack technique détaillée de VibraGuard}
\end{table}

\section{Capteurs IoT et Justification des Choix}

Le nœud de collecte de \projectname{} est basé sur la carte \textbf{ESP32} \cite{espressif2024}, microcontrôleur
puissant avec WiFi intégré, ADC 12 bits et deux cœurs de traitement, permettant de calculer
et d'envoyer les données sans module externe. Trois capteurs complémentaires sont utilisés :

\begin{table}[H]
  \centering
  \begin{tabularx}{\textwidth}{L{2.5cm} L{4cm} L{5.5cm}}
    \toprule
    \rowcolor{gray!30}
    \textbf{Capteur} &
    \textbf{Grandeur mesurée} &
    \textbf{Justification vs. alternative} \\
    \midrule
    \rowcolor{white}
    DS18B20 & Température (°C) & Interface numérique 1-Wire directe, pas de bruit de conversion.
    \textit{vs. thermocouple} : ce dernier exige un amplificateur externe. \\
    \rowcolor{gray!12}
    ACS712 & Courant RMS (A) et THD (\%) & Mesure sans rupture du circuit, isolation galvanique.
    \textit{vs. résistance shunt} : expose le microcontrôleur à la haute tension. \\
    \rowcolor{white}
    Piézo + ADS1115 & Vibration RMS, Peak, Kurtosis (mm/s) & ADC 16 bits (ADS1115) :
    4× plus de résolution que l'ADC interne ESP32. \textit{vs. ADC interne} :
    les vibrations faibles seraient noyées dans le bruit ADC 12 bits. \\
    \bottomrule
  \end{tabularx}
  \caption{Capteurs IoT du nœud VibraGuard et justification des choix}
\end{table}

\section{Structure des Données Collectées}

Chaque relevé du moteur (identifié par \texttt{motor\_id}) contient 6 indicateurs produits
par les trois capteurs :

\begin{table}[H]
  \centering
  \begin{tabularx}{\textwidth}{L{3cm} L{2.5cm} C{1.5cm} L{5cm}}
    \toprule
    \rowcolor{gray!30}
    \textbf{Donnée} &
    \textbf{Capteur} &
    \textbf{Unité} &
    \textbf{Ce qu'elle révèle} \\
    \midrule
    \rowcolor{white}       Température              & DS18B20           & °C      & Surcharge thermique, ventilation insuffisante \\
    \rowcolor{gray!12}Courant RMS              & ACS712            & A       & Charge mécanique, blocage, court-circuit \\
    \rowcolor{white}       THD Courant              & ACS712      & \%      & Harmoniques : qualité réseau, défauts électriques \\
    \rowcolor{gray!12}Vibration RMS            & Piézo + ADS1115   & mm/s    & Énergie globale, déséquilibre \\
    \rowcolor{white}       Vibration Peak           & Piézo + ADS1115   & mm/s    & Chocs ponctuels, impacts mécaniques \\
    \rowcolor{gray!12}Kurtosis                 & Piézo + ADS1115   & ---     & Détection précoce défauts roulements \\
    \bottomrule
  \end{tabularx}
  \caption{Indicateurs collectés par capteur — 6 champs par relevé}
\end{table}

\section{Formules de Calcul des Capteurs}

Le capteur DS18B20 transmet la température directement en degrés Celsius via le protocole
numérique 1-Wire (plage : $-55$°C à $+125$°C, précision $\pm 0{,}5$°C) \cite{ds18b20}. Aucune conversion
n'est nécessaire côté applicatif. Les seuils d'interprétation sont :

\begin{table}[H]
  \centering
  \begin{tabularx}{\textwidth}{C{3cm} L{8cm}}
    \toprule
    \rowcolor{gray!30}
    \textbf{Plage} & \textbf{Interprétation} \\
    \midrule
    \rowcolor{white}       0 -- 60\,°C   & Fonctionnement normal \\
    \rowcolor{gray!12}60 -- 75\,°C  & Alerte — charge élevée \\
    \rowcolor{white}       75 -- 90\,°C  & Surchauffe critique \\
    \rowcolor{gray!12}$>$ 90\,°C   & Risque de panne / arrêt machine \\
    \bottomrule
  \end{tabularx}
  \caption{Zones d'interprétation — Température DS18B20}
\end{table}

\subsection{ACS712 — Courant RMS (\texttt{current\_rms}) (A)}

Le capteur ACS712 exploite l'effet Hall : le courant crée un champ magnétique proportionnel,
converti en tension par le capteur. Le calcul du courant RMS se fait en quatre étapes sur
$N$ échantillons ADC 12 bits (0--4095) de l'ESP32 :

\begin{enumerate}
  \item \textbf{Conversion ADC → Tension :}
  \begin{equation}
    V = \frac{\text{ADC}}{4095} \times 3{,}3 \quad [\text{ADC} \in [0, 4095],\; V_{\max} = 3{,}3\,\text{V}]
  \end{equation}

  \item \textbf{Retrait de l'offset (sortie nulle sans courant = 2,5 V \cite{acs712}) :}
  \begin{equation}
    V_{\text{centré}} = V - 2{,}5
  \end{equation}

  \item \textbf{Calcul du courant instantané} (sensibilité ACS712-5A : $S = 0{,}185\,\text{V/A}$ \cite{acs712}) :
  \begin{equation}
    I = \frac{V_{\text{centré}}}{0{,}185}
  \end{equation}

  \item \textbf{Calcul RMS sur $N$ échantillons :}
  \begin{equation}
    I_{\text{RMS}} = \sqrt{\frac{I_1^2 + I_2^2 + \cdots + I_N^2}{N}}
  \end{equation}
\end{enumerate}

États du courant RMS : Normal $0{,}5$--$4$\,A | Charge élevée $4$--$5$\,A | Surcharge $> 5$\,A.

\subsection{ACS712 — THD Courant (\texttt{current\_thd}) (\%)}

Le Taux de Distorsion Harmonique (THD) mesure la qualité de la forme d'onde du courant.
Un THD élevé indique une distorsion harmonique pouvant provoquer une surchauffe du moteur.
Son calcul s'effectue via FFT sur $N$ échantillons :

\begin{enumerate}
  \item Acquérir $N$ échantillons du courant instantané $I(t)$.
  \item Appliquer la FFT (transformation signal temporel → fréquentiel).
  \item Identifier le fondamental $I_1$ (50\,Hz) et les harmoniques $I_2$ (100\,Hz), $I_3$ (150\,Hz)...
\end{enumerate}

\begin{equation}
  \text{THD} = \frac{\sqrt{I_2^2 + I_3^2 + \cdots + I_n^2}}{I_1} \times 100\%
\end{equation}

\noindent\textit{Exemple numérique :} $I_1 = 3{,}2\,\text{A}$, $I_2 = 0{,}5\,\text{A}$, $I_3 = 0{,}2\,\text{A}$ :
$\text{THD} = \frac{\sqrt{0{,}5^2 + 0{,}2^2}}{3{,}2} \times 100 = 16{,}8\%$

Seuils THD : Normal $0$--$5\%$ | Début distorsion $5$--$10\%$ | Distorsion visible $10$--$25\%$ | Danger $> 25\%$.

\subsection{Piézo + ADS1115 — Vibration RMS (\texttt{vib\_rms}) (mm/s)}

L'ADS1115 \cite{ads1115} (ADC 16 bits, plage $-32768$ à $+32767$, $V_{\text{ref}} = 4{,}096\,\text{V}$)
convertit le signal piézoélectrique avec une précision 4 fois supérieure à l'ADC interne de
l'ESP32. Le calcul du Vibration RMS suit les étapes :

\begin{enumerate}
  \item \textbf{Conversion brute → tension :}
  \begin{equation}
    V_i = \text{ADC} \times \frac{V_{\text{ref}}}{32767} = \text{ADC} \times \frac{4{,}096}{32767}
  \end{equation}

  \item \textbf{Retrait de l'offset (mesure au repos) :}
  \begin{equation}
    V_{\text{offset}} = \text{mean}(V_i) \quad \text{[moteur au repos]}
  \end{equation}

  \item \textbf{Calcul RMS et conversion en mm/s} (Sensibilité : $S = 20{,}0\,\text{mV/(mm/s)}$) selon le capteur et la calibration :
  \begin{equation}
    \text{Vib\_RMS} = \frac{\sqrt{\text{mean}((V_i - V_{\text{offset}})^2)} \times 1000}{S}
  \end{equation}
\end{enumerate}

États : Machine stable $0{,}5$--$3{,}0$\,mm/s | Début déséquilibre $3{,}0$--$5{,}0$ | Défaut mécanique $5{,}0$--$10{,}0$ | Risque panne $> 10{,}0$.

\subsection{Piézo + ADS1115 — Vibration Peak (\texttt{vib\_peak}) (mm/s)}

La Vibration Peak mesure la valeur maximale absolue enregistrée sur la période de mesure.
Elle permet de détecter les chocs ou impulsions soudaines :

\begin{multline}
  \text{Vib\_Peak} = \frac{\max(|V_i - V_{\text{offset}}|) \times 1000}{S} \\
  \text{(Sensibilité : } S = 20{,}0\,\text{mV/(mm/s) selon le capteur et la calibration)}
\end{multline}

États : Machine saine $< 3$\,mm/s | Début de défaut $3$--$6$ | Défaut mécanique clair $6$--$12$ | Risque panne grave $> 12$.

\subsection{Piézo + ADS1115 — Kurtosis (\texttt{vib\_kurtosis})}

Le Kurtosis est un indicateur statistique très sensible aux chocs impulsionnels \cite{antoni2006}. Il est
principalement utilisé pour la détection précoce des défauts de roulements. Une valeur $> 3$
signale des chocs anormaux :

\begin{equation}
  K = \frac{\text{mean}\left((V_i - \mu)^4\right)}{\left(\text{mean}\left((V_i - \mu)^2\right)\right)^2}
  \quad \text{avec } \mu = \text{mean}(V_i)
\end{equation}

\noindent\textit{Exemple :} $V_1=0{,}2$\,V, $V_2=0{,}5$\,V, $V_3=0{,}4$\,V, $\mu = 0{,}367$\,V :
\[
K = \frac{\frac{(0{,}2-\mu)^4 + (0{,}5-\mu)^4 + (0{,}4-\mu)^4}{3}}
        {\left(\frac{(0{,}2-\mu)^2 + (0{,}5-\mu)^2 + (0{,}4-\mu)^2}{3}\right)^2}
\]

Interprétation : $K \approx 3$ (signal gaussien normal) | $K > 3$ (chocs, début défaut roulement) | $K < 3$ (signal plat, peu de variation).

Zones : Normal $2$--$4$ | Légère anomalie $4$--$6$ | Vibrations non linéaires $6$--$10$ | Chocs forts / défaut roulement $> 10$.


\subsection{Simulation des Envois MQTT — Script \texttt{external\_mqtt\_sender.py}}

Dans le cadre du développement et des tests, un script de simulation MQTT a été développé
afin de reproduire le comportement des capteurs IoT réels. La figure ci-dessous montre
l'exécution de ce script, qui génère et envoie des trames JSON vers le broker MQTT de
\projectname{} (\texttt{vibraguard.mywire.org:30083}) :

\begin{figure}[H]
  \centering
  \includegraphics[width=\textwidth]{l_envoi_des_requetes_mqtt_semulees.png}
  \caption{Exécution du script \texttt{external\_mqtt\_sender.py} --- Simulation d'envois
  MQTT vers le broker \texttt{vibraguard.mywire.org} avec données de capteurs en temps réel}
\end{figure}

Chaque trame MQTT publiée sur le topic \texttt{vibraguard/sensors} contient les champs
\texttt{motor\_id}, \texttt{vib\_rms}, \texttt{vib\_peak},
\texttt{vib\_kurtosis}, \texttt{current\_rms}, \texttt{current\_thd} et
\texttt{temperature}. On observe ici des données correspondant au moteur
\texttt{MTR-Broyeur-04} avec un statut \texttt{anomalous}, permettant de valider la
chaîne complète de détection d'anomalies en conditions simulées avant le déploiement
sur les équipements réels de l'OCP Safi.

\section{Formules de Calcul des KPI}

Les indicateurs de performance clés (KPI) affichés sur le dashboard de \projectname{} sont
calculés selon les formules industrielles standardisées suivantes :

\subsection{MTBF — Mean Time Between Failures}

Le \textbf{MTBF} mesure le temps moyen de fonctionnement entre deux pannes successives.
Un MTBF élevé indique une bonne fiabilité de l'équipement :

\begin{equation}
  \text{MTBF} = \frac{\text{Temps Total de Fonctionnement}}{\text{Nombre de Pannes}}
  = \frac{\sum_{i=1}^{n} T_{\text{fonctionnement},i}}{n}
\end{equation}

\subsection{MTTR — Mean Time To Repair}

Le \textbf{MTTR} mesure le temps moyen nécessaire pour réparer un équipement après une
défaillance. Un MTTR faible indique une maintenance efficace :

\begin{equation}
  \text{MTTR} = \frac{\text{Temps Total de Réparation}}{\text{Nombre de Pannes}}
  = \frac{\sum_{i=1}^{n} T_{\text{réparation},i}}{n}
\end{equation}

\subsection{Taux de Disponibilité}

Le \textbf{taux de disponibilité} (ou disponibilité opérationnelle) exprime la proportion
du temps pendant laquelle un équipement est opérationnel :

\begin{equation}
  D = \frac{\text{MTBF}}{\text{MTBF} + \text{MTTR}} \times 100\%
\end{equation}

\subsection{Précision du Modèle Random Forest}

La précision globale du classifieur est évaluée par :

\begin{equation}
  \text{Accuracy} = \frac{\text{VP} + \text{VN}}{\text{VP} + \text{VN} + \text{FP} + \text{FN}}
\end{equation}

Le F1-Score par classe, plus robuste en cas de déséquilibre de classes, est donné par :

\begin{equation}
  F_1 = 2 \times \frac{\text{Précision} \times \text{Rappel}}{\text{Précision} + \text{Rappel}}
  = \frac{2 \cdot \text{VP}}{2 \cdot \text{VP} + \text{FP} + \text{FN}}
\end{equation}

\section{Conclusion}

Ce chapitre a décrit en détail l'architecture microservices de \projectname{}, du nœud
de collecte IoT jusqu'aux dashboards décisionnels. L'architecture retenue garantit la
scalabilité horizontale, la résilience et l'observabilité complète du système. Les capteurs
ESP32 (DS18B20, ACS712, Piézo+ADS1115) offrent une couverture multi-dimensionnelle des
paramètres moteurs, tandis que le pipeline Kafka--Spark assure un traitement haute performance.
Le déploiement sur GCP Kubernetes avec CI/CD automatisé constitue le socle opérationnel de
la solution. Le chapitre suivant présente la mise en œuvre effective du pipeline de données
et les résultats du modèle d'intelligence artificielle.

% ============================
\chapter{Implémentation du Pipeline de Données et de l'IA}
% ============================

Ce chapitre décrit le pipeline de données de \projectname{}, de l'acquisition des données capteurs réelles à l'inférence en temps réel, en passant par le prétraitement et l'entraînement du modèle de Machine Learning.

\section{Pipeline de Prétraitement et Entraînement}

\subsection{Structure des Données d'Entraînement}

Le dataset d'entraînement (\texttt{sensor\_data\_training.csv}) a été exploré et validé
via le script \texttt{explore\_data.py}. La figure suivante présente la sortie de cette
exploration, confirmant la structure et la qualité des données :

\begin{figure}[H]
  \centering
  \includegraphics[width=\textwidth]{structure_de_data_pour_le_modele_ia.png}
  \caption{Exploration du dataset d'entraînement (\texttt{sensor\_data\_training.csv}) ---
  Structure des 11 colonnes, aperçu des 5 premières lignes et résumé des types de données}
\end{figure}

\begin{infobox}[Précision sur l'origine des données]
Le dataset utilisé pour l'entraînement du modèle Random Forest a été fourni par le
\textbf{Département IT d'OCP Safi}. Ces données ont été collectées par des
étudiants de la filière Systèmes Embarqués dans le cadre d'un \textbf{environnement
de simulation} reproduisant les plages de valeurs et les signatures vibratoires,
thermiques et électriques caractéristiques des états moteurs identifiés.

Cette approche par simulation a permis de pallier l'absence d'historique de pannes
labellisées suffisant pour certaines classes de défauts rares, tout en respectant
des plages physiques réalistes et cohérentes avec les capteurs (ESP32, DS18B20,
ACS712, Piézo+ADS1115).
\end{infobox}



\subsection{Limites liées à la Taille du Dataset}
Le volume de 1,000 échantillons constitue une limite méthodologique qu'il convient d'expliciter. Ce choix s'explique par le fait que les données ont été \textbf{collectées dans un environnement de simulation} par des étudiants de la filière Systèmes Embarqués, insuffisante pour constituer un corpus massif couvrant l'ensemble des classes de défauts avec un grand nombre d'occurrences — notamment pour les classes de défauts rares (déséquilibre, désalignement, surcharge électrique).
La séparation du dataset en ensembles d’entraînement (80%) et de test (20%) a été réalisée avec l’option stratify=y de la fonction train_test_split de Scikit-learn. Cette stratification garantit que les proportions de chaque classe sont préservées dans les deux sous-ensembles. La validation croisée ($  k=5  $) a par ailleurs été utilisée pour vérifier la stabilité du score de précision.

Néanmoins, la généralisation du modèle sur un parc industriel élargi (perspective de déploiement pilote sur 5 à 10 moteurs réels de l’OCP Safi) nécessitera un réentraînement sur un corpus substantiellement plus large et issu de mesures réelles de terrain, conformément à la recommandation de réentraînement périodique formulée en conclusion.

Le dataset contient \textbf{1 000 échantillons} répartis sur \textbf{11 colonnes} :
\texttt{timestamp}, \texttt{sample\_id}, \texttt{moteur\_id},
\texttt{vib\_rms}, \texttt{vib\_peak}, \texttt{vib\_kurtosis}, \texttt{current\_rms},
\texttt{current\_thd}, \texttt{temperature}, \texttt{label} et \texttt{anomaly\_type}.
On note la présence de valeurs manquantes sur certaines colonnes numériques (\texttt{vib\_rms} et \texttt{vib\_peak} : 969/1000, \texttt{current\_rms} :
976/1000, \texttt{temperature} : 967/1000), justifiant l'étape d'imputation par la médiane
dans le pipeline de prétraitement. Les colonnes \texttt{label} et \texttt{anomaly\_type}
sont de type \texttt{object} (catégoriel) et constituent les cibles de classification.

\subsection{Étapes de Prétraitement}

Le pipeline de prétraitement est implémenté avec la bibliothèque Scikit-learn \cite{sklearn},
qui fournit des transformateurs composables via l'API \texttt{Pipeline} :

\begin{enumerate}
  \item \textbf{Imputation :} remplacement des valeurs manquantes par la médiane
  (\texttt{SimpleImputer}, \texttt{strategy='median'}), méthode recommandée pour les
  données de capteurs à distribution asymétrique \cite{sklearn}.
  \item \textbf{Filtrage (Clipping) :} limitation des valeurs aberrantes à $\pm 3\sigma$
  pour l'exclusion d'outliers.
  \item \textbf{Normalisation :} mise à l'échelle par \texttt{StandardScaler} pour
  garantir la convergence du modèle.
  \item \textbf{Séparation train/test :} split 80/20 stratifié, préservant les proportions
  de classes dans les deux sous-ensembles.
\end{enumerate}

\subsection{Modélisation — Random Forest Classifier}

Le Random Forest \cite{breiman2001} a été choisi pour sa robustesse face au surapprentissage, sa capacité à gérer les données de haute dimension et sa bonne interprétabilité via l'importance des variables. 

Il a été entraîné avec les hyperparamètres suivants :

\begin{itemize}
  \item \texttt{n\_estimators = 100} : Ce nombre d'arbres offre un bon compromis entre la stabilité des prédictions et le temps de calcul. Un nombre trop faible augmente la variance, tandis qu'un nombre trop élevé augmente inutilement le temps d'entraînement sans gain significatif en performance.
  
  \item \texttt{max\_depth = 10} : Cette profondeur limite le surapprentissage tout en permettant au modèle de capturer des relations complexes dans les données de vibrations industrielles.
  
  \item \texttt{min\_samples\_split = 5} : Cette contrainte permet d'éviter la création de nœuds trop spécifiques, améliorant ainsi la généralisation du modèle sur de nouvelles données.
\end{itemize}

La figure ci-dessous illustre l'exécution complète du script d'entraînement
\texttt{train\_random\_forest.py}, depuis le chargement des données jusqu'à la
sérialisation des modèles :

\begin{figure}[H]
  \centering
  \includegraphics[width=\textwidth]{l_entrainement_de_modele_ia.png}
  \caption{Exécution du script \texttt{train\_random\_forest.py} --- Entraînement du
  modèle Random Forest avec une précision globale élevée et sérialisation \texttt{joblib}}
\end{figure}

La sortie console confirme le déroulement correct du pipeline : chargement de 1 000 samples,
nettoyage des données (imputation et filtrage outliers), entraînement du classifieur
Random Forest, puis sauvegarde des artefacts \texttt{vibraguard\_rf\_model.joblib} et
\texttt{vibraguard\_scaler.joblib}. La précision globale mesurée lors de cet entraînement
atteint \textbf{92,00\%}.

\begin{resultbox}
  \textbf{Résultat clé :} Le modèle Random Forest \cite{breiman2001} atteint une précision globale
  de \textbf{92,0\%} sur six classes d'état, dépassant largement l'objectif de 85\% fixé dans
  le cahier des charges. Ce niveau de performance est cohérent avec les résultats rapportés par
  Lei et al. \cite{lei2020} pour des applications similaires de diagnostic de machines tournantes.
  Les modèles entraînés (scaler + classifier) sont sérialisés via \texttt{joblib} pour déploiement
  dans le pipeline Spark Streaming \cite{spark}.
\end{resultbox}


\subsection{Séparation Stratifiée Train/Test et Matrice de Confusion}

La séparation du dataset en ensembles d'entraînement (80\%) et de test (20\%) a été
réalisée avec l'option \textbf{\texttt{stratify=y}} de la fonction
\texttt{train\_test\_split} de Scikit-learn \cite{sklearn}. Cette option garantit que
les proportions de chaque classe sont \textbf{préservées à l'identique} dans les deux
sous-ensembles, évitant ainsi tout biais d'échantillonnage — particulièrement critique
dans notre contexte où la classe \texttt{NONE} (état normal) est majoritaire avec
environ 64\% des échantillons, tandis que les classes de défauts (\texttt{DEFAUT\_ROULEMENT},
\texttt{DESALIGNEMENT}, \texttt{DESEQUILIBRE}, \texttt{SURCHARGE\_ELECTRIQUE},
\texttt{SURCHAUFFE}) sont minoritaires. Sans \texttt{stratify=y}, le tirage aléatoire
pourrait produire un ensemble de test sous-représentant certaines classes de défauts,
rendant l'évaluation non représentative des conditions réelles d'exploitation à l'OCP Safi.

\begin{figure}[H]
  \centering
  \includegraphics[width=0.85\textwidth]{confusion_matrix.png}
  \caption{Matrice de confusion du modèle Random Forest sur l'ensemble de test ---
  Classification des six états moteurs avec une précision globale de 92,0\%}
  \label{fig:confusion_matrix}
\end{figure}

La matrice de confusion (figure~\ref{fig:confusion_matrix}) confirme les excellentes
performances du modèle Random Forest sur l'ensemble de test. L'analyse par classe
révèle les résultats suivants :

\begin{itemize}
  \item \textbf{NONE (état normal) :} 128 prédictions correctes sur 128 échantillons,
  soit un rappel de \textbf{100\%}. Aucune confusion avec les autres classes, ce qui
  est fondamental pour éviter les fausses alarmes en production.
  \item \textbf{DESALIGNEMENT :} 12 vrais positifs sur 13, avec une seule confusion
  vers la classe \texttt{NONE}.
  \item \textbf{DEFAUT\_ROULEMENT :} 11 vrais positifs sur 14, avec 3 confusions
  vers \texttt{NONE}.
  \item \textbf{DESEQUILIBRE :} 9 vrais positifs sur 11, avec 2 confusions vers
  \texttt{NONE}.
  \item \textbf{SURCHARGE\_ELECTRIQUE :} 12 vrais positifs sur 16, avec 4 confusions
  vers \texttt{NONE}.
  \item \textbf{SURCHAUFFE :} 12 vrais positifs sur 18, avec 1 confusion vers
  \texttt{DEFAUT\_ROULEMENT} et 3 vers \texttt{DESEQUILIBRE} et 2 vers \texttt{NONE}.
\end{itemize}



\subsection{Rapport de Classification Détaillé}

Au-delà de la précision globale et de la matrice de confusion, un \textbf{rapport de
classification} (\textit{classification report}), généré via la fonction
\texttt{classification\_report} de Scikit-learn \cite{sklearn}, fournit une évaluation
fine de la performance du modèle pour chacune des six classes. Trois métriques
complémentaires sont calculées par classe :

\begin{itemize}
  \item \textbf{Précision (\textit{precision}) 92.56\% :} proportion de prédictions positives
  correctes parmi toutes les prédictions positives pour une classe donnée. Une précision
  élevée pour une classe de défaut signifie un faible taux de fausses alarmes.
  
  \item \textbf{Rappel (\textit{recall}) 92.00\% :} proportion d'échantillons réellement
  positifs correctement détectés. Un rappel élevé est critique pour les classes de
  défauts, car un défaut non détecté (faux négatif) peut entraîner une panne réelle.
  
  \item \textbf{Score F1 (\textit{F1-score}) 92.00\% :} moyenne harmonique de la précision et
  du rappel, particulièrement pertinente pour les classes minoritaires en cas de
  déséquilibre des classes.
\end{itemize}


\section{Traitement en Streaming avec Apache Spark}

Le script \texttt{spark\_streaming\_process.py} constitue le moteur de traitement temps réel,
exploitant les capacités de micro-batch processing d'Apache Spark Streaming \cite{spark} :

\begin{enumerate}
  \item Connexion au cluster Kafka \cite{kafka} et abonnement aux topics de données capteurs.
  \item Réception des micro-batchs (fenêtre de 2 secondes).
  \item Désérialisation JSON et validation du schéma.
  \item Application du pipeline de prétraitement (scaler \texttt{joblib} \cite{sklearn}).
  \item Inférence par le modèle Random Forest \cite{breiman2001} : classification de l'état.
  \item Mise à jour des KPI (score de santé, compteur d'anomalies, MTBF estimé) conformément
  aux indicateurs de performance définis par la norme ISO 13374 \cite{iso13374}.
  \item Publication des résultats vers le Backend via REST API \cite{springboot}.
  \item Déclenchement d'alertes pour les anomalies détectées.
\end{enumerate}

\subsection{Diagramme de Séquence : Détection d'Anomalie → Alerte (Création Manuelle d'OT)}

\begin{figure}[H]
  \centering
  \includegraphics[width=\textwidth]{sequence_anomalie_ot.png}
  \caption{Diagramme de séquence : Flux de détection d'anomalie jusqu'à la création manuelle d'un Ordre de Travail}
  \label{fig:seq-anomalie-ot}
\end{figure}

\section{Conclusion}

Ce chapitre a démontré la robustesse du pipeline de données de \projectname{}, depuis
l'acquisition des signaux capteurs jusqu'à l'inférence en temps réel. Le modèle Random
Forest atteint une précision de \textbf{92,0\%} sur six classes de défauts, dépassant
significativement l'objectif initial de 85\%. Le pipeline de prétraitement (imputation,
filtrage $\pm 3\sigma$, normalisation) garantit la qualité des données d'entrée, tandis que
le traitement Spark Streaming assure une latence de traitement d'environ 1,2 seconde. Ces
résultats valident la viabilité de l'approche IoT--Big Data--IA pour la maintenance
prédictive. Le chapitre suivant présente les composants applicatifs développés (backend,
frontend et intégration Blockchain).

% ============================
\chapter{Développement du Backend, Frontend et Blockchain}
% ============================

Ce chapitre présente le processus de développement de \projectname{} par phases, les composants applicatifs réalisés (backend microservices, interface web, Blockchain/IPFS) et la présentation de l'application finale.

\section{Processus de Développement}

Le développement de \projectname{} est organisé en \textbf{cinq phases} articulées selon
l'approche Agile, chacune correspondant à un ensemble de tâches Jira.

\subsection{Phase 1 : Préparation}

\begin{itemize}
  \item Configuration de l'environnement DevOps : dépôt GitHub, branches \texttt{main}/\texttt{develop}/\texttt{feature/*}, pipeline CI/CD GitHub Actions.
  \item Mise en place de l'infrastructure pour le développement local (Kafka, Mosquitto, Oracle, MongoDB, IPFS) via Kubernetes.
  \item Installation et paramétrage du réseau blockchain local avec Hardhat \cite{hardhat} (testnet privé).
  \item Définition des conventions de codage, des standards API (OpenAPI 3.0) et des templates de PR GitHub.
\end{itemize}

\subsection{Phase 2 : Conception et Design}

\begin{itemize}
  \item Modélisation UML complète : diagrammes de cas d'utilisation, classes et séquence.
  \item Conception de l'architecture microservices (diagramme C4 : contexte, conteneurs,
  composants, code) selon les principes de Spring Boot \cite{springboot}.
  \item Design des interfaces utilisateurs (wireframes Figma) : dashboard, vue moteur, gestion OT,
  en accord avec les principes d'ergonomie des applications industrielles.
  \item Spécification des APIs REST (Swagger/OpenAPI 3.0).
  \item Conception du schéma de base de données (Oracle + MongoDB).
  \item Design des Smart Contracts : interfaces Solidity \cite{ethereum}, événements, fonctions,
  gestion des accès (RBAC on-chain).
\end{itemize}

\subsection{Phase 3 : Développement Backend}

\begin{itemize}
  \item Implémentation des microservices Spring Boot \cite{springboot} avec injection de dépendances (IoC).
  \item Intégration du pipeline Kafka \cite{kafka} (Producers/Consumers) et Spark Streaming \cite{spark} (jobs PySpark).
  \item Développement des repositories JPA/Hibernate pour Oracle DB et Spring Data MongoDB pour les vibrations.
  \item Implémentation du Smart Contract Solidity \cite{ethereum} (\texttt{WorkOrderRegistry.sol})
  déployé via Hardhat \cite{hardhat}.
  \item Déploiement du nœud IPFS \cite{ipfs} et intégration via HTTP API Spring Boot \cite{springboot}.
  \item Mise en place du monitoring (Prometheus metrics), conformément aux
  exigences de disponibilité ISO 27001 \cite{iso27001}.
\end{itemize}

\subsection{Phase 4 : Développement Frontend}

\begin{itemize}
  \item \textbf{Application Web (React 18 \cite{react} + TypeScript) :} composants réutilisables
  (MotorCard, AlertBadge, KPIWidget), intégration WebSocket pour le temps réel.
  \item Intégration des appels REST vers le backend avec gestion des tokens JWT, conformément
  aux recommandations de sécurité ISO 27001 \cite{iso27001}.
  \item Tests unitaires (Jest pour React \cite{react}) et tests E2E (Cypress).
\end{itemize}

\subsection{Phase 5 : Tests et Déploiement}

\begin{itemize}
  \item \textbf{Tests d'intégration :} couverture 78\% (objectif $>$70\%) testant les
  interactions inter-microservices et les flux Kafka \cite{kafka} end-to-end.
  \item \textbf{Tests de charge :} simulation de 10 000 capteurs simultanés avec Apache
  JMeter --- résultats : débit 12 000 msg/s, latence P95 = 145 ms.
  \item \textbf{Tests de sécurité :} scan OWASP ZAP sur le backend, audit du Smart
  Contract \cite{ethereum} (revue manuelle), test de pénétration basique selon ISO 27001 \cite{iso27001}.
  \item \textbf{Déploiement Kubernetes \cite{kubernetes} :} manifests Helm Charts,
  auto-scaling HPA (CPU $>$70\%), health probes (liveness/readiness), rolling updates sans downtime.
  \item \textbf{Documentation :} Swagger UI généré automatiquement, guides utilisateur
  (technicien et responsable maintenance), rapport PFE final.
\end{itemize}

La figure suivante présente la sortie de la commande \texttt{kubectl get all -n vibraguard},
confirmant l'état opérationnel de l'ensemble des pods, services et déploiements dans le
namespace \projectname{} sur le cluster Kubernetes :

\begin{figure}[H]
  \centering
  \includegraphics[width=\textwidth]{les_pods_kubernetes.png}
  \caption{État des ressources Kubernetes dans le namespace \texttt{vibraguard} ---
  Pods, services, déploiements et replica sets tous en statut \texttt{Running}/\texttt{Ready}}
\end{figure}

On constate que l'ensemble des \textbf{pods} sont opérationnels (\texttt{Running}), attestant de la stabilité de la plateforme. Les
microservices déployés incluent : \texttt{backend}, \texttt{blockchain}, \texttt{frontend},
\texttt{ipfs}, \texttt{kafka}, \texttt{mosquitto}, \texttt{mqtt-kafka-bridge},
\texttt{oracle-db}, \texttt{mongodb} et \texttt{spark}. Les services exposent leurs ports
respectifs (\texttt{backend:30007}, \texttt{frontend:30008}, \texttt{mosquitto:30083},
\texttt{oracle-db:31255}), et l'ensemble des déploiements affichent un ratio
\texttt{READY 1/1}, confirmant le déploiement sans downtime réalisé via rolling update.

\section{Architecture Backend — Microservices Spring Boot}

\subsection{API Gateway}

L'API Gateway (Spring Cloud Gateway) constitue le point d'entrée unique avec :
\begin{itemize}
  \item Routage dynamique vers les microservices selon le path de la requête.
  \item Authentification centralisée via JWT.
  \item Rate limiting pour la protection contre les abus.
  \item Logging centralisé de toutes les requêtes.
\end{itemize}

\subsection{Endpoints REST API}

Les endpoints REST de \projectname{} sont organisés par domaine fonctionnel, selon les
contrôleurs Spring Boot (\texttt{@RestController}) implémentés :

\begin{table}[H]
  \centering
  \scriptsize
  \begin{tabularx}{\textwidth}{L{5.5cm} C{1.8cm} L{5cm}}
    \toprule
    \rowcolor{gray!30}
    \textbf{Endpoint} & \textbf{Méthode} & \textbf{Description} \\
    \midrule
    \rowcolor{white}       \texttt{/api/v1/auth/register} & POST & Création d'un nouveau compte utilisateur \\
    \rowcolor{gray!12}\texttt{/api/v1/auth/login}    & POST & Authentification et génération du token JWT \\
    \rowcolor{white}       \texttt{/api/v1/auth/me}       & GET  & Profil de l'utilisateur connecté \\
    \bottomrule
  \end{tabularx}
  \caption{Endpoints REST --- Authentification (\texttt{/api/v1/auth})}
\end{table}

\begin{table}[H]
  \centering
  \scriptsize
  \begin{tabularx}{\textwidth}{L{5.5cm} C{1.8cm} L{5cm}}
    \toprule
    \rowcolor{gray!30}
    \textbf{Endpoint} & \textbf{Méthode} & \textbf{Description} \\
    \midrule
    \rowcolor{white}       \texttt{/api/v1/iot/motors}                    & GET    & Liste tous les moteurs avec état de santé \\
    \rowcolor{gray!12}\texttt{/api/v1/iot/motors}                    & POST   & Ajoute un nouveau moteur \\
    \rowcolor{white}       \texttt{/api/v1/iot/motors/\{id\}}             & GET    & Détails d'un moteur spécifique \\
    \rowcolor{gray!12}\texttt{/api/v1/iot/motors/\{id\}}             & PUT    & Mise à jour d'un moteur \\
    \rowcolor{white}       \texttt{/api/v1/iot/motors/\{id\}}             & DELETE & Supprime le moteur et ses données \\
    \rowcolor{gray!12}\texttt{/api/v1/iot/motors/\{id\}/vibration}   & GET    & Historique des vibrations d'un moteur \\
    \rowcolor{white}       \texttt{/api/v1/iot/motors/vibrations}         & GET    & Toutes les données de vibration \\
    \rowcolor{gray!12}\texttt{/api/v1/iot/motors/vibrations}         & POST   & Enregistre une nouvelle lecture capteur \\
    \bottomrule
  \end{tabularx}
  \caption{Endpoints REST --- Moteurs \& Vibrations (\texttt{/api/v1/iot/motors})}
\end{table}

\begin{table}[H]
  \centering
  \scriptsize
  \begin{tabularx}{\textwidth}{L{5.5cm} C{1.8cm} L{5cm}}
    \toprule
    \rowcolor{gray!30}
    \textbf{Endpoint} & \textbf{Méthode} & \textbf{Description} \\
    \midrule
    \rowcolor{white}       \texttt{/api/v1/ml/alerts}               & GET  & Liste les alertes (filtrées par rôle) \\
    \rowcolor{gray!12}\texttt{/api/v1/ml/alerts}               & POST & Crée une nouvelle alerte \\
    \rowcolor{white}       \texttt{/api/v1/ml/alerts/\{id\}}        & PUT  & Met à jour statut, priorité, niveau \\
    \rowcolor{gray!12}\texttt{/api/v1/ml/alerts/mark-all-read} & POST & Marque toutes les alertes comme lues \\
    \rowcolor{white}       \texttt{/api/v1/ml/alerts/\{id\}/read}   & POST & Marque une alerte spécifique comme lue \\
    \bottomrule
  \end{tabularx}
  \caption{Endpoints REST --- Alertes ML (\texttt{/api/v1/ml/alerts})}
\end{table}

\begin{table}[H]
  \centering
  \scriptsize
  \begin{tabularx}{\textwidth}{L{5.5cm} C{1.8cm} L{5cm}}
    \toprule
    \rowcolor{gray!30}
    \textbf{Endpoint} & \textbf{Méthode} & \textbf{Description} \\
    \midrule
    \rowcolor{white}       \texttt{/api/v1/iot/work-orders}        & GET    & Liste les OT (filtrés par rôle) \\
    \rowcolor{gray!12}\texttt{/api/v1/iot/work-orders}        & POST   & Crée un OT et notifie le technicien \\
    \rowcolor{white}       \texttt{/api/v1/iot/work-orders/\{id\}} & PUT    & Met à jour statut, assignation, coût \\
    \rowcolor{gray!12}\texttt{/api/v1/iot/work-orders/\{id\}} & DELETE & Supprime un ordre de travail \\
    \bottomrule
  \end{tabularx}
  \caption{Endpoints REST --- Ordres de Travail (\texttt{/api/v1/iot/work-orders})}
\end{table}

\begin{table}[H]
  \centering
  \scriptsize
  \begin{tabularx}{\textwidth}{L{5.5cm} C{1.8cm} L{5cm}}
    \toprule
    \rowcolor{gray!30}
    \textbf{Endpoint} & \textbf{Méthode} & \textbf{Description} \\
    \midrule
    \rowcolor{white}       {\tiny\texttt{/api/v1/iot/technicians}}          & GET    & Liste tous les techniciens \\
    \rowcolor{gray!12}{\tiny\texttt{/api/v1/iot/technicians/\{id\}}}   & GET    & Profil d'un technicien \\
    \rowcolor{white}       {\tiny\texttt{/api/v1/iot/technicians/\{id\}}}   & PUT    & Mise à jour du profil \\
    \rowcolor{gray!12}{\tiny\texttt{/api/v1/iot/technicians/\{id\}}}   & DELETE & Supprime un compte technicien \\
    \rowcolor{white}       {\tiny\texttt{/api/v1/bi/kpis}}                  & GET    & KPI temps réel (MTTR, disponibilité, coûts) \\
    \rowcolor{gray!12}{\tiny\texttt{/api/v1/bi/mtbf-by-site}}          & GET    & MTBF regroupé par site \\
    \rowcolor{white}       {\tiny\texttt{/api/v1/bi/maintenance-costs}}     & GET    & Historique et distribution des coûts \\
    \rowcolor{gray!12}{\tiny\texttt{/api/v1/bi/interventions}}         & GET    & Statistiques sur les interventions \\
    \bottomrule
  \end{tabularx}
  \caption{Endpoints REST --- Techniciens \& BI ({\tiny\texttt{/api/v1/iot/technicians}}, \texttt{/api/v1/bi})}
\end{table}

\section{Interface Utilisateur Web — React}

Le frontend web est développé en React 18 avec TypeScript et Tailwind CSS. Il comprend :

\subsection{Dashboard Principal}
\begin{itemize}
  \item \textbf{KPI globaux :} taux de disponibilité, score de santé moyen, alertes actives, OT en cours.
  \item \textbf{Liste des moteurs} avec codes couleur adaptatifs (Vert / Jaune / Orange / Rouge).
  \item \textbf{Graphiques temps réel} des mesures vibratoires (WebSocket pour les mises à jour push).
  \item \textbf{Fil d'alertes récentes} avec priorité et type de défaut.
\end{itemize}

\subsection{Vue Détaillée d'un Moteur}
\begin{itemize}
  \item \textbf{SanteCard :} jauge circulaire représentant le score de santé (0--100\%).
  \item Graphiques de tendance des paramètres vibratoires, thermiques et électriques.
  \item Historique des anomalies et des ordres de travail associés.
\end{itemize}

\section{Intégration Blockchain et IPFS}

\subsection{Smart Contracts (Solidity)}

Le Smart Contract déployé sur le réseau Hardhat \cite{hardhat} est écrit en Solidity \cite{ethereum},
le langage de programmation natif de la blockchain Ethereum \cite{buterin2014} :
\begin{itemize}
  \item \texttt{WorkOrderRegistry.sol :} enregistrement immuable des ordres de travail
  avec les fonctions \texttt{createWorkOrder()}, \texttt{updateStatus()}, \texttt{verifyIntegrity()}.
  Cette approche garantit une auditabilité conforme aux exigences de la norme ISO 55000 \cite{iso55000}.
\end{itemize}

\subsection{Intégration IPFS}

L'intégration IPFS \cite{ipfs} repose sur le protocole de stockage décentralisé développé par
Protocol Labs, qui garantit l'adressage par contenu (\textit{content-addressed storage}) :

\begin{enumerate}
  \item Le rapport PDF est généré.
  \item Le fichier est uploadé vers IPFS \cite{ipfs} via l'API HTTP.
  \item IPFS retourne un CID (Content Identifier) --- hash cryptographique unique du fichier,
  calculé selon l'algorithme SHA2-256 \cite{ipfs}.
  \item L'immutabilité est garantie : toute modification ultérieure du fichier génèrerait un CID
  différent, rendant toute falsification détectable.
\end{enumerate}

\subsection{Diagramme de Séquence : Génération Rapport PDF → Certification Blockchain/IPFS}

\begin{figure}[H]
  \centering
  \includegraphics[width=0.85\textwidth]{sequence_rapport_blockchain.png}
  \caption{Diagramme de séquence : Génération et certification immuable d’un rapport via IPFS + Blockchain}
  \label{fig:seq-rapport-blockchain}
\end{figure}


\section{Business Intelligence — Tableaux de Bord BI}

Des tableaux de bord BI offrent une dimension analytique et décisionnelle :
\begin{itemize}
  \item \textbf{Dashboard Opérationnel :} suivi temps réel des alertes et OT.
  \item \textbf{Dashboard Analytique :} tendances MTBF/MTTR, analyse coûts des pannes, performance par équipement.
\end{itemize}

\section{Présentation de l'Application}

\subsection{Architecture Générale de l'Interface}

L'interface utilisateur de \projectname{} est organisée autour de trois espaces distincts
selon le profil de l'utilisateur connecté :

\begin{itemize}
  \item \textbf{Espace Technicien :} centré sur les alertes actives et les ordres de travail
  assignés. Comprend : liste OT, formulaire d'intervention, consultation diagnostic moteur.
  \item \textbf{Espace Responsable Maintenance :} centré sur la supervision globale du parc.
  Dashboard web React avec KPI temps réel (MTBF, MTTR, disponibilité), carte thermique des
  moteurs, historiques et analyses de tendances.
  \item \textbf{Espace Administrateur :} gestion des utilisateurs (RBAC), configuration des
  moteurs.
\end{itemize}

\subsection{Fonctionnalités Clés de l'Application Web}

\begin{itemize}
  \item \textbf{Dashboard Temps Réel :} actualisation automatique via
  WebSocket. Affiche les KPI globaux, la carte du parc de moteurs avec codes couleur adaptatifs
  (Vert = Normal, Orange = Attention, Rouge = Critique) et le fil des dernières alertes.

  \item \textbf{Vue Détaillée Moteur :} page dédiée à chaque moteur avec jauge de santé
  circulaire (score 0--100\%), graphiques de tendance multi-paramètres (vibration,
  température, courant), timeline des anomalies et liste des OT associés.

  \item \textbf{Gestion des Ordres de Travail :} tableau Kanban des OT (\`{A} faire / En cours /
  Termin\'{e}), formulaire de création lié à une anomalie, assignation technicien.

  \item \textbf{Module Blockchain :} page de vérification d'intégrité permettant de valider la conformité des données avec l'enregistrement on-chain.

  \item \textbf{Dashboard BI embarqués :} onglets analytiques avec les tableaux de bord
  KPI, tendances MTBF/MTTR et analyse des coûts de maintenance,
  génération PDF du rapport avec enregistrement automatique IPFS.
\end{itemize}

\subsection{Sécurité de l'Application}

La sécurité de \projectname{} est assurée à plusieurs niveaux, conformément aux exigences
de la norme ISO/IEC 27001 \cite{iso27001} en matière de sécurité des systèmes d'information :

\begin{itemize}
  \item \textbf{Authentification :} JWT avec refresh tokens (expiration 15 min) pour toutes
  les requêtes vers l'API Gateway \cite{springboot}.
  \item \textbf{Autorisation :} RBAC (\textit{Role-Based Access Control}) granulaire avec
  3 rôles (Technicien, Responsable, Administrateur), tant au niveau applicatif Spring \cite{springboot}
  que contractuel (Solidity \cite{ethereum}).
  \item \textbf{Audit :} toutes les actions utilisateurs sont tracées dans les logs, garantissant une traçabilité
  inaltérable et auditée, en ligne avec les recommandations ISO 55000 \cite{iso55000}.
\end{itemize}


% ============================================================
\section{Présentation Illustrée de l'Application}
% ============================================================

Cette section présente les principales captures d'écran de la plateforme \projectname{},
illustrant le processus de développement (gestion de projet, API, bases de données) et
les interfaces utilisateurs livrées.

% ---- Sous-section : Outils de développement ----
\subsection{Outils et Environnement de Développement}

\subsubsection{Tableau de Bord Jira — Gestion Agile du Projet}

\begin{figure}[H]
  \centering
  \includegraphics[width=0.85\textwidth]{1.jira_screencapture.png}
  \caption{Tableau Kanban Jira du projet VibraGuard --- Vue des colonnes \og{}À faire\fg{},
  \og{}En cours\fg{}, \og{}Terminé\fg{} et \og{}Deploy\fg{} avec les tickets}
  \label{fig:jira}
\end{figure}

Le tableau Kanban Jira (figure~\ref{fig:jira}) illustre la gestion Agile du projet.
On distingue les quatre colonnes d'état : \textbf{À faire}, \textbf{En cours},
\textbf{Terminé} et \textbf{Deploy}.

\subsubsection{Documentation API — Swagger UI}

\begin{figure}[H]
  \centering
  \includegraphics[width=0.85\textwidth]{2.swagger_screencapture.png}
  \caption{Documentation Swagger UI de l'API VibraGuard --- Vue de l'ensemble des
  endpoints REST organisés par domaine fonctionnel (Authentication, Motors, Alerts,
  WorkOrders, Technicians, Analytics)}
  \label{fig:swagger}
\end{figure}

La figure~\ref{fig:swagger} présente la documentation interactive Swagger UI générée
automatiquement à partir des annotations OpenAPI 3.0 des contrôleurs Spring Boot
\cite{springboot}. Les six groupes d'endpoints — \textit{Authentication},
\textit{Motors}, \textit{Alerts}, \textit{WorkOrders}, \textit{Technicians} et
\textit{Analytics} — reflètent directement l'architecture microservices de
\projectname{} et les tableaux d'endpoints REST détaillés dans la section précédente.
Les schémas de requête/réponse (\texttt{RegisterRequest}, \texttt{Motor},
\texttt{VibrationData}, \texttt{WorkOrder}, \texttt{KPIData}, etc.) sont également
exposés, facilitant l'intégration avec des systèmes tiers.

\subsubsection{Base de Données Oracle — Structure des Tables Métier}

\begin{figure}[H]
  \centering
  \includegraphics[width=0.75\textwidth]{3.oracle_Capture.png}
  \caption{Connexion Oracle Database 21c Express Edition et structure des tables
  \texttt{USERS} et \texttt{MOTORS} via SQL*Plus dans le namespace Kubernetes de
  \projectname{}}
  \label{fig:oracle}
\end{figure}

La figure~\ref{fig:oracle} confirme le déploiement opérationnel d'Oracle Database 21c
Express Edition dans le namespace Kubernetes \texttt{vibraguard}. La commande
\texttt{DESC USERS} révèle les colonnes de la table des utilisateurs (\texttt{ID},
\texttt{EMAIL}, \texttt{FULL\_NAME}, \texttt{PASSWORD}, \texttt{ROLE}, \texttt{STATUS},
\texttt{DEPARTMENT}, \texttt{ENTERPRISE}, \texttt{EMPLOYEE\_ID},
\texttt{PHONE\_NUMBER}), modélisant le contrôle d'accès RBAC de la plateforme. La
commande \texttt{DESC MOTORS} expose la structure des moteurs (\texttt{ID},
\texttt{ETAT\_PCT}, \texttt{ETAT\_LABEL}, \texttt{ETAT\_COLOR},
\texttt{LOCALISATION}, \texttt{DERNIERE\_ALERTE\_TYPE}, \texttt{POWER},
\texttt{SITE}), entité centrale du domaine métier de \projectname{}.

\subsubsection{Base de Données MongoDB — Collection des Vibrations}

\begin{figure}[H]
  \centering
  \includegraphics[width=0.85\textwidth]{4.mongodb_Capture.png}
  \caption{Connexion au pod MongoDB 8.2.9 via \texttt{mongosh} dans Kubernetes ---
  Base de données \texttt{vibraguard} et collection \texttt{vibration\_data} opérationnelles}
  \label{fig:mongodb}
\end{figure}

La figure~\ref{fig:mongodb} atteste du fonctionnement de MongoDB 8.2.9 déployé dans
le namespace Kubernetes \texttt{vibraguard}. La commande \texttt{show dbs} confirme
l'existence de la base \texttt{vibraguard} (40 KiB), et \texttt{show collections}
expose la collection \texttt{vibration\_data}, réceptacle des séries temporelles
vibratoires générées par le pipeline Kafka--Spark. MongoDB a été retenu pour ce stockage
en raison de sa flexibilité schéma-less, particulièrement adaptée aux données IoT dont
la structure peut évoluer avec l'ajout de nouveaux capteurs.

% ---- Sous-section : Interfaces utilisateurs ----
\subsection{Interfaces Utilisateurs de la Plateforme Web}

\subsubsection{Page d'Inscription}

\begin{figure}[H]
  \centering
  \includegraphics[width=0.82\textwidth]{5.register_image17.png}
  \caption{Page d'inscription de \projectname{} --- Formulaire de création de compte
  avec sélection du rôle (Administrateur, Responsable, Technicien) et de l'entreprise}
  \label{fig:register}
\end{figure}

La page d'inscription (figure~\ref{fig:register}) permet la création d'un compte sur
la plateforme \projectname{}. Le formulaire collecte le nom complet, l'adresse
électronique professionnelle, le mot de passe (avec confirmation et indicateur de
force), l'entreprise et le rôle (\textit{Administrateur}, \textit{Responsable
Maintenance} ou \textit{Technicien}). L'enregistrement déclenche une requête
\texttt{POST /api/v1/auth/register}, qui crée l'utilisateur en base Oracle DB et
retourne un token JWT permettant l'accès immédiat à la plateforme.

\subsubsection{Page de Connexion}

\begin{figure}[H]
  \centering
  \includegraphics[width=0.75\textwidth]{6.login_image18.png}
  \caption{Page de connexion sécurisée de \projectname{} --- Authentification JWT
  avec option \og{}Se souvenir de moi\fg{} et lien vers la création de compte}
  \label{fig:login}
\end{figure}

La page de connexion (figure~\ref{fig:login}) constitue le point d'entrée sécurisé de
la plateforme. L'utilisateur saisit son adresse électronique professionnelle et son mot
de passe ; une requête \texttt{POST /api/v1/auth/login} est soumise au backend Spring
Boot \cite{springboot}, qui valide les identifiants et retourne un token JWT (durée de
vie 15 minutes) avec un refresh token. L'option \og{}Se souvenir de moi\fg{} prolonge
la session via le mécanisme de refresh token, conformément aux bonnes pratiques de
sécurité ISO/IEC 27001 \cite{iso27001}.

\subsubsection{Tableau de Bord Principal}

\begin{figure}[H]
  \centering
  \includegraphics[width=\textwidth]{7.home_screencapture.png}
  \caption{Tableau de bord principal de \projectname{} --- KPI temps réel, cartographie
  santé des moteurs, graphique d'analyse vibratoire et fil des alertes récentes}
  \label{fig:dashboard}
\end{figure}

Le tableau de bord principal (figure~\ref{fig:dashboard}) constitue la vue centrale de
\projectname{} pour le profil Administrateur/Responsable. Il agrège quatre KPI globaux
en temps réel : le nombre total de moteurs surveillés (2), le nombre de moteurs en état
critique (1), le taux de disponibilité global (30,0\%) et le nombre d'alertes actives
(7). La \textbf{cartographie santé} présente visuellement les moteurs MTR-Broyeur-04 et
MTR-Broyeur-07 avec leur statut colorimétrique. Le \textbf{graphique d'analyse Vibration
RMS} affiche les données temps réel reçues via WebSocket. La section \og{}Moteurs Sous
Surveillance\fg{} détaille par moteur l'état de santé (30\,\% Critique pour MTR-Broyeur-04
vs. 100\,\% Normal pour MTR-Broyeur-07) et la vibration RMS courante. Enfin, le
\textbf{fil des alertes récentes} liste les dernières anomalies détectées par le modèle
Random Forest, avec le type de défaut et la valeur vibratoire associée.

\subsubsection{Liste des Moteurs}

\begin{figure}[H]
  \centering
  \includegraphics[width=\textwidth]{8.motors_screencapture.png}
  \caption{Page de gestion des moteurs industriels --- Vue carte des deux moteurs
  MTR-Broyeur-04 et MTR-Broyeur-07 avec leur état de santé et vibration RMS courante}
  \label{fig:motors}
\end{figure}

La page des moteurs (figure~\ref{fig:motors}) offre une vue inventaire du parc de
machines surveillées. Chaque carte moteur affiche l'identifiant (\texttt{MTR-Broyeur-04},
\texttt{MTR-Broyeur-07}), la localisation (Safi / Zone 1, Secteur A / Niveau 1),
l'état de santé (\og{}Optimal\fg{} en vert au moment de la capture, indiquant l'absence
d'anomalie active), la vibration RMS instantanée (1,2 mm/s) et le type de moteur
(Asynchrone). Le bouton \og{}Ajouter un Moteur\fg{} permet l'enregistrement d'un
nouvel équipement via l'endpoint \texttt{POST /api/v1/iot/motors}. Les filtres de
recherche (par zone et par état de santé) facilitent la supervision d'un parc étendu.

\subsubsection{Centre des Alertes}

\begin{figure}[H]
  \centering
  \includegraphics[width=\textwidth]{9.alertes_screencapture.png}
  \caption{Page des alertes de \projectname{} --- Liste des 7 alertes critiques actives
  classées par sévérité, avec type de défaut, score de confiance IA et statut}
  \label{fig:alertes}
\end{figure}

La page des alertes (figure~\ref{fig:alertes}) centralise l'ensemble des anomalies
détectées par le modèle Random Forest \cite{breiman2001} en temps réel. Les 7 alertes
actives concernent le moteur MTR-Broyeur-04 et couvrent trois types de défauts :
\texttt{DEFAUT\_ROULEMENT} (score de confiance IA : 53,77\%), \texttt{DESALIGNEMENT}
(48,4\,\% à 81,92\,\%) et \texttt{SURCHARGE\_ELECTRIQUE} (47\,\%). Toutes sont
classifiées \og{}Critique\fg{} par le système. Les filtres de sévérité, de plage de
dates et de moteur permettent une navigation efficace dans l'historique des alertes.
Les boutons d'action autorisent la confirmation ou l'escalade de chaque alerte.

\subsubsection{Gestion des Ordres de Travail (Kanban)}

\begin{figure}[H]
  \centering
  \includegraphics[width=\textwidth]{11.ot_screencapture.png}
  \caption{Vue Kanban des Ordres de Travail de \projectname{} --- Colonnes \og{}À faire\fg{},
  \og{}En cours\fg{} et \og{}Terminé\fg{} avec les OT W-458 et W-459}
  \label{fig:ot_kanban}
\end{figure}

La vue Kanban des Ordres de Travail (figure~\ref{fig:ot_kanban}) reprend le paradigme
de gestion Agile appliqué à la maintenance industrielle. L'OT \textbf{W-458} (priorité
Haute, moteur MTR-Broyeur-04, type Correctif, durée estimée 1h27, coût 400 MAD, assigné
à Youssef Bouali) est en colonne \og{}À faire\fg{}. L'OT \textbf{W-459} (priorité Haute,
moteur MTR-Broyeur-07, durée 2h30, coût 800 MAD) est \og{}En cours\fg{}. Le bouton
\og{}Nouveau OT\fg{} déclenche le formulaire de création d'ordre de travail.

\subsubsection{Formulaire de Création d'un Ordre de Travail}

\begin{figure}[H]
  \centering
  \includegraphics[width=\textwidth]{12.add_ot_image22.png}
  \caption{Formulaire de création d'un Ordre de Travail dans \projectname{} ---
  Sélection du moteur, du type d'anomalie, de la date, du technicien assigné et des
  pièces nécessaires}
  \label{fig:add_ot}
\end{figure}

Le formulaire de création d'OT (figure~\ref{fig:add_ot}) permet au responsable de
maintenance de formaliser une intervention. Les champs requis incluent : le moteur
concerné (liste déroulante liée à l'API \texttt{GET /api/v1/iot/motors}), le type
d'anomalie et la sévérité associée, la date planifiée et la durée estimée (en heures
et minutes), le type de maintenance (\og{}Préventif\fg{} ou \og{}Correctif\fg{}), le
technicien assigné (liste déroulante des techniciens disponibles), les pièces
nécessaires (recherche dans le magasin) et le coût estimé en MAD. La soumission du
formulaire déclenche un \texttt{POST /api/v1/iot/work-orders} et notifie automatiquement
le technicien assigné.

\subsubsection{Audit Blockchain}

\begin{figure}[H]
  \centering
  \includegraphics[width=\textwidth]{13.audit_blockchain_image26.png}
  \caption{Page d'Audit Blockchain de \projectname{} --- Tableau de bord d'intégrité
  avec hash de transaction, numéro de bloc, moteur concerné et bouton de vérification}
  \label{fig:audit_blockchain}
\end{figure}

La page d'audit Blockchain (figure~\ref{fig:audit_blockchain}) expose les KPI
d'intégrité de la chaîne : 1 bloc sécurisé, 1 contrat intelligent actif sur le réseau
Hardhat interne, un taux d'intégrité de \textbf{100\%} (aucune altération détectée) et
un temps de synthèse (validation) de 0 seconde. La section \og{}Traçabilité
d'Intervention en Direct\fg{} liste les transactions enregistrées on-chain, avec leur
hash complet, le
numéro de bloc (\#2), le moteur concerné (MTR-Broyeur-04), l'intervention (\og{}Création
OT : Sans titre\fg{}), la date et l'adresse Ethereum de l'utilisateur.
Le bouton \og{}Vérifier intégrité\fg{} déclenche l'appel au Smart Contract
\texttt{verifyIntegrity()} \cite{ethereum,hardhat}.

\subsubsection{Rapports IPFS}

\begin{figure}[H]
  \centering
  \includegraphics[width=\textwidth]{14.ipfs_image25.png}
  \caption{Page Rapports IPFS de \projectname{} --- Rapport PDF quotidien stocké sur
  IPFS avec son CID
  et compteur de téléchargements}
  \label{fig:ipfs_reports}
\end{figure}

La page Rapports IPFS (figure~\ref{fig:ipfs_reports}) liste les rapports PDF générés
par la plateforme et stockés de façon décentralisée. Le rapport est stocké sur IPFS. Son CID constitue son empreinte cryptographique unique, générée à partir de son contenu \cite{ipfs} : toute modification ultérieure du fichier générerait
un CID différent, trahissant immédiatement toute tentative de falsification. Le bouton
\og{}Générer Rapport\fg{} déclenche la génération PDF côté serveur, son upload vers
le nœud IPFS.

\subsubsection{Tableau de Bord Décisionnel BI}

\begin{figure}[H]
  \centering
  \includegraphics[width=\textwidth]{15.bi_image24.png}
  \caption{Tableau de Bord Décisionnel BI de \projectname{} --- KPI MTBF/MTTR/Disponibilité,
  analyse des coûts de maintenance, cartographie géographique des sites (Maroc)
  et répartition des types d'intervention}
  \label{fig:bi}
\end{figure}

Le Tableau de Bord Décisionnel BI (figure~\ref{fig:bi}) fournit aux responsables
maintenance une vue analytique agrégée. Les KPI du mois d'avril 2026 affichent :
un MTBF quotidien de \textbf{24,0\,h}, un MTTR de \textbf{4,5\,h} (calculs conformes
aux formules ISO 13374 \cite{iso13374} détaillées au Chapitre III), une disponibilité
de \textbf{50\,\%} et un coût de maintenance (période test).
Le graphique de coûts de maintenance compare le coût réel au budget alloué sur la
période. Le diagramme en anneau de répartition des interventions indique 50\,\%
d'interventions préventives. La \textbf{cartographie géographique du Maroc} localise
les sites OCP connectés, affichant 1 site connecté et 55 alertes actives. Le module
est accessible en vue quotidienne, hebdomadaire ou mensuelle, avec export PDF et Excel.


\section{Conclusion}

Ce chapitre a présenté l'ensemble des composants applicatifs de \projectname{}, depuis
les microservices Spring Boot jusqu'à l'interface utilisateur React, en passant par
l'intégration Blockchain/IPFS et dashboard BI. Le processus de développement
Agile en cinq phases (préparation, conception, backend, frontend, tests/déploiement) a
permis de livrer une plateforme complète et opérationnelle. La sécurité à plusieurs niveaux
(JWT, RBAC, audit trail Blockchain) garantit l'intégrité et la confidentialité des données
de maintenance. Les tests d'intégration et de charge
confirment la solidité de l'architecture. La conclusion générale synthétise les résultats
obtenus et trace les perspectives d'évolution de \projectname{}.

% ============================
\setcounter{chapter}{0}
\chapter*{Conclusion Générale}
\addcontentsline{toc}{chapter}{Conclusion Générale}
\markboth{Conclusion Générale}{Conclusion Générale}
% ============================

Le projet \projectname{} a permis de concevoir, développer et valider une plateforme complète
de maintenance prédictive pour les moteurs industriels, répondant pleinement aux objectifs fixés.
Sur le plan du traitement massif de données, l'architecture combinant Apache Kafka \cite{kafka}
et Apache Spark \cite{spark} s'est révélée pleinement opérationnelle, le pipeline ayant démontré
sa capacité à absorber un débit élevé de messages tout en maintenant une
latence de traitement faible. Sur le volet intelligence artificielle, le modèle
Random Forest \cite{breiman2001} entraîné dans le cadre du projet atteint une précision globale
de 92,0\% sur six classes de défauts moteurs, un résultat qui dépasse l'objectif initial
et qui se révèle cohérent avec l'état de l'art rapporté pour des applications comparables de
diagnostic de machines tournantes \cite{lei2020}. L'architecture applicative repose quant à elle
sur un ensemble de microservices Spring Boot \cite{springboot} déployés sur Kubernetes
\cite{kubernetes}, bénéficiant d'un auto-scaling natif et d'une chaîne d'intégration et de
déploiement continus entièrement automatisée. La dimension Blockchain et stockage décentralisé
n'a pas été négligée : l'intégration d'Ethereum et d'IPFS \cite{ethereum, ipfs} permet de certifier
de façon immuable l'ensemble des rapports, avec un audit trail complet garantissant
la conformité aux normes,
l'application web développée en React \cite{react} offre une expérience utilisateur moderne et
réactive, restituant les données en temps réel grâce à une couche d'observabilité.

Le tableau ci-après synthétise, pour chacun des critères de succès retenus en début de projet,
l'objectif fixé et le résultat effectivement obtenu. Il en ressort que tous les critères ont été
atteints, et que la majorité d'entre eux --- précision de détection, débit de traitement, ainsi
que la traçabilité blockchain --- ont même été dépassés, ce qui confirme la pertinence des choix
architecturaux et technologiques opérés tout au long du projet.

\begin{table}[H]
  \centering
  \begin{tabularx}{\textwidth}{L{4.5cm} C{2.5cm} C{2.5cm} C{2cm}}
    \toprule
    \rowcolor{gray!30}
    \textbf{Critère} &
    \textbf{Objectif} &
    \textbf{Résultat} &
    \textbf{Statut} \\
    \midrule
    \rowcolor{white}       Précision détection  & $> 85\%$         & $92{,}0\%$          & \textbf{Dépassé} \\
    \rowcolor{gray!12}Latence traitement   & $< 2$\,s         & $\approx 1{,}2$\,s  & \textbf{Atteint} \\
    \rowcolor{white}       Débit de traitement  & $> 10\,000$\,msg/s & $> 12\,000$\,msg/s    & \textbf{Dépassé} \\
    \rowcolor{gray!12}Temps réponse API    & $< 200$\,ms (P95)& $\approx 145$\,ms   & \textbf{Atteint} \\
    \rowcolor{white}       Disponibilité        & $> 99{,}5\%$     & $99{,}7\%$          & \textbf{Atteint} \\
    \rowcolor{white}       Traçabilité blockchain & $100\%$        & $100\%$             & \textbf{Atteint} \\
    \rowcolor{gray!12}Couverture tests     & $> 70\%$         & $78\%$              & \textbf{Atteint} \\
    \bottomrule
  \end{tabularx}
  \caption{Critères de succès — Objectifs vs. Résultats atteints}
\end{table}

\section*{Perspectives d'Amélioration}

Si \projectname{} répond pleinement aux objectifs fixés dans le cadre de ce Projet de Fin
d'Études, plusieurs axes d'amélioration se dessinent pour les mois et les années à venir,
selon un horizon temporel croissant.

À court terme, deux pistes complémentaires méritent d'être explorées. La première concerne
l'apprentissage profond, et plus particulièrement les réseaux de neurones récurrents de type
LSTM \cite{hochreiter1997}, dont l'aptitude à modéliser des dépendances temporelles longues
pourrait enrichir l'analyse des séries vibratoires sur des périodes étendues, en complément du
Random Forest actuellement déployé. La seconde piste relève de l'explicabilité des modèles
(\textit{Explainable AI}) : l'intégration d'outils tels que SHAP (SHapley Additive exPlanations) pour une explication globale des prédictions et LIME (Local Interpretable Model-agnostic Explanations) pour une interprétation locale et compréhensible de chaque prédiction permet de fournir aux techniciens non spécialistes des explications claires et compréhensibles, renforçant ainsi la confiance et l'adoption opérationnelle de la solution \cite{lei2020}.

À moyen terme, l'évolution de \projectname{} pourrait s'orienter vers l'\textit{Edge Computing}
et le \textit{TinyML}, en déployant des modèles allégés directement sur les microcontrôleurs IoT
afin de réduire la latence et de permettre une détection d'anomalies même en l'absence de
connectivité réseau. Le développement d'un jumeau numérique (\textit{Digital Twin}) pour chaque
moteur constituerait une seconde avancée significative, en autorisant la simulation de scénarios
de dégradation avant qu'ils ne surviennent réellement sur le terrain. Enfin, une migration vers
Hyperledger Fabric pourrait être envisagée afin de bénéficier de meilleures performances
transactionnelles et d'une gouvernance d'entreprise mieux adaptée à un déploiement à grande échelle.

À plus long terme, trois perspectives structurantes se dégagent. Le développement d'une
application mobile dotée de capacités de réalité augmentée permettrait aux techniciens de
visualiser directement les données de diagnostic sur l'équipement physique, en superposition de
la réalité. L'intégration de \projectname{} au système ERP du Groupe OCP automatiserait quant à
elle les commandes de pièces de rechange ainsi que l'imputation comptable des coûts de
maintenance. Enfin, l'adoption de techniques d'apprentissage fédéré (\textit{Federated Learning})
ouvrirait la voie à une amélioration collaborative des modèles de détection entre plusieurs
sites industriels du Groupe, sans qu'aucun d'entre eux n'ait à partager ses données brutes,
préservant ainsi la confidentialité tout en mutualisant les bénéfices de l'apprentissage.

% ============================================================
\appendix
\inappendixparttrue
% ============================================================

\chapter{Glossaire des Termes Techniques}

\begin{longtable}{L{3.5cm} L{9cm}}
  \toprule
  \rowcolor{gray!30}
  \textbf{Terme} & \textbf{Définition} \\
  \midrule
  \endhead
  \bottomrule
  \endfoot
  \rowcolor{white}       API              & Interface de programmation permettant la communication entre applications. \\
  \rowcolor{white}       CID              & Content Identifier : identifiant cryptographique unique d'un fichier IPFS. \\
  \rowcolor{gray!12}DAP              & Diammonium Phosphate : engrais phosphaté produit par réaction de l'acide phosphorique avec l'ammoniac, utilisé en agriculture. \\
  \rowcolor{white}       ETL              & Extract Transform Load : processus d'extraction, transformation et chargement de données. \\
  \rowcolor{gray!12}IIoT             & Industrial Internet of Things : IoT appliqué aux systèmes industriels. \\
  \rowcolor{white}       IPFS             & InterPlanetary File System : protocole de stockage décentralisé pair-à-pair à adressage par contenu développé par Protocol Labs. \\
  \rowcolor{gray!12}JWT              & JSON Web Token : standard ouvert pour l'échange sécurisé d'informations. \\
  \rowcolor{white}       KPI              & Key Performance Indicator : indicateur clé de performance. \\
  \rowcolor{gray!12}MAP              & Monoammonium Phosphate : engrais phosphaté soluble, utilisé en agriculture et comme retardateur de flamme. \\
  \rowcolor{white}       MTBF             & Mean Time Between Failures : temps moyen entre deux pannes successives. \\
  \rowcolor{gray!12}MTTR             & Mean Time To Repair : temps moyen de réparation après une panne. \\
  \rowcolor{white}       MQTT             & Message Queuing Telemetry Transport : protocole léger de messagerie IoT standardisé par l'OASIS. \\
  \rowcolor{gray!12}OASIS            & Organization for the Advancement of Structured Information Standards : organisme de normalisation des standards informatiques, notamment MQTT. \\
  \rowcolor{white}       RBAC             & Role-Based Access Control : contrôle d'accès basé sur les rôles. \\
  \rowcolor{gray!12}SLA              & Service Level Agreement : accord de niveau de service définissant les engagements de qualité et de disponibilité d'un service. \\
  \rowcolor{white}       Smart Contract   & Programme autonome auto-exécutable déployé sur la blockchain. \\
  \rowcolor{gray!12}TSP              & Triple Superphosphate : engrais phosphaté concentré obtenu par traitement de la roche phosphatée à l'acide phosphorique. \\
\end{longtable}

\chapter{Travaux Effectués et Résultats}

\section{Récapitulatif des Travaux Effectués}

L'ensemble des travaux réalisés dans le cadre du projet \projectname{} peut être regroupé
en cinq grands domaines d'activité :

\begin{table}[H]
  \centering
  \small
  \begin{tabularx}{\textwidth}{L{3cm} L{8cm} C{1.5cm}}
    \toprule
    \rowcolor{gray!30}
    \textbf{Domaine} &
    \textbf{Travaux Réalisés} &
    \textbf{Statut} \\
    \midrule
    \rowcolor{white}       Data Science    & Pipeline ETL, entraînement Random Forest (92,0\%), prétraitement et évaluation & \textbf{Done} \\
    \rowcolor{gray!12}Big Data        & Configuration cluster Kafka, jobs Spark Streaming (latence 1,2s), stockage MongoDB & \textbf{Done} \\
    \rowcolor{white}       Backend         & Microservices Spring Boot, APIs REST, sécurité JWT, RBAC & \textbf{Done} \\
    \rowcolor{gray!12}Frontend        & Dashboard React temps réel (WebSocket), gestion OT, alertes & \textbf{Done} \\
    \rowcolor{white}       Blockchain/IPFS & Smart Contract Solidity (WorkOrderRegistry), intégration IPFS, audit trail & \textbf{Done} \\
    \rowcolor{gray!12}Infrastructure  & Déploiement Kubernetes GCP, CI/CD GitHub Actions, monitoring Prometheus/Grafana & \textbf{Done} \\
    \bottomrule
  \end{tabularx}
  \caption{Récapitulatif des travaux effectués}
\end{table}

\section{Résultats et Recommandations}

\subsection{Résultats Techniques}

\begin{resultbox}
  \textbf{Synthèse des résultats obtenus :}\\
  Précision du modèle IA : \textbf{92,0\%} (objectif : 85\%) |
  Latence : \textbf{1,2 s} (objectif : $<$ 2 s) |
  Disponibilité : \textbf{99,7\%} (objectif : 99,5\%) |
  Débit : \textbf{12 000 msg/s} (objectif : 10 000 msg/s)
\end{resultbox}

\subsection{Analyse des Résultats}

\begin{itemize}
  \item \textbf{Précision IA :} le dépassement de 10 points de l'objectif initial (92,0\%
  vs. 85\%) s'explique par la qualité des données capteurs réelles, la robustesse du pipeline de
  nettoyage et le choix pertinent du Random Forest.

  \item \textbf{Performance Big Data :} la latence de 1,2 secondes est obtenue grâce à
  l'optimisation des micro-batchs Spark (fenêtres de 2s) et au tuning des paramètres du
  cluster Kafka (partitions, réplication, compression LZ4).

  \item \textbf{Fiabilité Blockchain :} 100\% des événements critiques sont certifiés
  on-chain, avec un temps de confirmation moyen de 3 secondes sur le réseau Hardhat privé,
  sans aucune transaction échouée lors des tests de charge.
\end{itemize}

\subsection{Recommandations}

À l'issue de ce projet, plusieurs recommandations sont formulées pour le déploiement en
production et la pérennisation de la solution :

\begin{enumerate}
  \item \textbf{Déploiement progressif :} commencer par un pilote sur 5 à 10 moteurs
  critiques de l'usine de Safi avant généralisation à l'ensemble du parc.
  \item \textbf{Formation des équipes :} organiser des sessions de formation pour les
  techniciens et responsables maintenance sur l'utilisation de l'interface web.
  \item \textbf{Intégration ERP :} connecter \projectname{} au système ERP de l'OCP pour
  automatiser les commandes de pièces de rechange et l'intégration comptable des coûts.
  \item \textbf{Réentraînement périodique :} mettre en place un pipeline MLOps pour le
  réentraînement mensuel automatique des modèles sur les nouvelles données de production.
  \item \textbf{Migration Hyperledger :} envisager une migration vers Hyperledger Fabric
  pour les besoins d'une blockchain enterprise avec gouvernance fine.
  \item \textbf{Audit de sécurité :} conduire un audit de sécurité (pentest) avant le
  déploiement en production, en ciblant notamment les Smart Contracts et l'API Gateway.
\end{enumerate}



% ---- Bibliographie ----
%\bibliographystyle{ieeetr}
\inmainchaptersfalse

\begin{thebibliography}{99}

\bibitem{ocpgroup}
  Groupe OCP,
  \textit{Rapport Annuel OCP S.A.},
  \url{https://www.ocpgroup.ma}, consulté en 2025.

\bibitem{usgs2025}
  U.S. Geological Survey (USGS),
  \textit{Mineral Commodity Summaries -- Phosphate Rock},
  U.S. Department of the Interior, January 2025.

\bibitem{schwab2016}
  K. Schwab,
  \textit{The Fourth Industrial Revolution},
  World Economic Forum, Geneva, 2016.

\bibitem{hermann2016}
  M. Hermann, T. Pentek, B. Otto,
  \textit{Design Principles for Industrie 4.0 Scenarios},
  49th Hawaii International Conference on System Sciences (HICSS), pp. 3928--3937, 2016.

\bibitem{mobley2002}
  R. K. Mobley,
  \textit{An Introduction to Predictive Maintenance},
  2nd ed., Butterworth-Heinemann, 2002.

\bibitem{benhanifia2025}
  A. Benhanifia, Z. Ben Cheikh, P. Moura Oliveira, A. Valente, J. Lima,
  \textit{Systematic review of predictive maintenance practices in the manufacturing sector},
  Intelligent Systems with Applications, vol. 26, p. 200501, 2025.

\bibitem{randall2011}
  R. B. Randall,
  \textit{Vibration-based Condition Monitoring: Industrial, Aerospace and Automotive Applications},
  John Wiley \& Sons, 2011.

\bibitem{dewell1984}
  D. L. Dewell and L. D. Mitchell,
  ``Detection of a Misaligned Disk Coupling Using Spectrum Analysis,''
  \textit{ASME Journal of Vibration, Acoustics, Stress, and Reliability in Design},
  vol. 106, no. 1, pp. 9--16, 1984.

\bibitem{randallandantoni2011}
  R. B. Randall and J. Antoni,
  ``Rolling element bearing diagnostics---A tutorial,''
  \textit{Mechanical Systems and Signal Processing},
  vol. 25, no. 2, pp. 485--520, 2011.

\bibitem{lee2014}
  J. Lee, F. Wu, W. Zhao, M. Ghaffari, L. Liao, D. Siegel,
  \textit{Prognostics and health management design for rotary machinery systems -- Reviews, methodology and applications},
  Mechanical Systems and Signal Processing, vol. 42, no. 1--2, pp. 314--334, 2014.

\bibitem{mqtt}
  OASIS Standard,
  \textit{MQTT Version 3.1.1 Specification},
  \url{https://docs.oasis-open.org/mqtt/mqtt/v3.1.1/mqtt-v3.1.1.html}, 2014.

\bibitem{kreps2011}
  J. Kreps, N. Narkhede, J. Rao,
  \textit{Kafka: a Distributed Messaging System for Log Processing},
  Proceedings of the NetDB Workshop, VLDB, 2011.

\bibitem{kafka}
  Apache Software Foundation,
  \textit{Apache Kafka Documentation},
  \url{https://kafka.apache.org/documentation/}, 2024.

\bibitem{spark}
  Apache Software Foundation,
  \textit{Apache Spark Documentation},
  \url{https://spark.apache.org/docs/latest/}, 2024.

\bibitem{breiman2001}
  L. Breiman,
  \textit{Random forests},
  Machine Learning, vol. 45, no. 1, pp. 5--32, 2001.

\bibitem{sklearn}
  F. Pedregosa et al.,
  \textit{Scikit-learn: Machine Learning in Python},
  Journal of Machine Learning Research, vol. 12, pp. 2825--2830, 2011.

\bibitem{nakamoto2008}
  S. Nakamoto,
  \textit{Bitcoin: A Peer-to-Peer Electronic Cash System},
  \url{https://bitcoin.org/bitcoin.pdf}, 2008.

\bibitem{buterin2014}
  V. Buterin,
  \textit{A Next-Generation Smart Contract and Decentralized Application Platform},
  Ethereum White Paper, 2014.

\bibitem{iso55000}
  ISO,
  \textit{ISO 55000:2014 -- Asset management -- Overview, principles and terminology},
  International Organization for Standardization, 2014.

\bibitem{iso13374}
  ISO,
  \textit{ISO 13374-1:2003 -- Condition monitoring and diagnostics of machines},
  International Organization for Standardization, 2003.

\bibitem{hardhat}
  Nomic Foundation,
  \textit{Hardhat -- Ethereum Development Environment},
  \url{https://hardhat.org/docs}, 2024.

\bibitem{ipfs}
  Protocol Labs,
  \textit{IPFS Documentation},
  \url{https://docs.ipfs.tech/}, 2024.

\bibitem{jardine2006}
  A. K. S. Jardine, D. Lin, D. Banjevic,
  \textit{A review on machinery diagnostics and prognostics implementing condition-based maintenance},
  Mechanical Systems and Signal Processing, vol. 20, no. 7, pp. 1483--1510, 2006.

\bibitem{itil2019}
  AXELOS,
  \textit{ITIL Foundation: ITIL 4 Edition},
  TSO (The Stationery Office), 2019.

\bibitem{springboot}
  Broadcom / VMware (Spring),
  \textit{Spring Boot Reference Documentation},
  \url{https://docs.spring.io/spring-boot/docs/current/reference/html/}, 2024.

\bibitem{ethereum}
  Ethereum Foundation,
  \textit{Ethereum Developer Documentation},
  \url{https://ethereum.org/en/developers/docs/}, 2024.

\bibitem{kubernetes}
  The Linux Foundation,
  \textit{Kubernetes Documentation},
  \url{https://kubernetes.io/docs/home/}, 2024.

\bibitem{react}
  Meta Open Source,
  \textit{React -- A JavaScript Library for Building User Interfaces},
  \url{https://react.dev/}, 2024.

\bibitem{espressif2024}
  Espressif Systems,
  \textit{ESP32 Series Datasheet}, Version 5.2,
  \url{https://www.espressif.com/sites/default/files/documentation/esp32_datasheet_en.pdf}, 2024.

\bibitem{ds18b20}
  Analog Devices (Maxim Integrated),
  \textit{DS18B20 Programmable Resolution 1-Wire Digital Thermometer},
  \url{https://www.analog.com/media/en/technical-documentation/data-sheets/ds18b20.pdf}.

\bibitem{acs712}
  Allegro MicroSystems,
  \textit{ACS712: Fully Integrated, Hall Effect-Based Linear Current Sensor},
  \url{https://www.allegromicro.com/en/products/sense/current-sensor-ics/integrated-current-sensors/acs712}.

\bibitem{ads1115}
  Texas Instruments,
  \textit{ADS111x 16-Bit ADC Datasheet}, SBAS444D,
  \url{https://www.ti.com/lit/ds/symlink/ads1115.pdf}, 2018.

\bibitem{doe2010}
  U.S. Department of Energy,
  \textit{Operations \& Maintenance Best Practices: A Guide to Achieving Operational Efficiency},
  Federal Energy Management Program, August 2010.

\bibitem{antoni2006}
  J. Antoni,
  ``The spectral kurtosis: a useful tool for characterising non-stationary signals,''
  \textit{Mechanical Systems and Signal Processing},
  vol.~20, no.~2, pp.~282--307, 2006.

\bibitem{lei2020}
  Y. Lei, B. Yang, X. Jiang, F. Jia, N. Li, A. K. Nandi,
  \textit{Applications of machine learning to machine fault diagnosis: A review and roadmap},
  Mechanical Systems and Signal Processing, vol. 138, 2020.

\bibitem{iso27001}
  ISO/IEC,
  \textit{ISO/IEC 27001:2022 -- Information security management systems},
  International Organization for Standardization, 2022.

\bibitem{hochreiter1997}
  S. Hochreiter, J. Schmidhuber,
  \textit{Long short-term memory},
  Neural Computation, vol. 9, no. 8, pp. 1735--1780, 1997.

\end{thebibliography}

\end{document}